
%********************************************
\section{Rotor disk theory}
%*******************************************

%**************************
\subsection{Introduction}
%**************************

Disk rotor theory is often refered as  "Froude theory".

It is the basic one-dimensional inviscid theory first developped to analyse
mass flow, pressure drop, momentun, thrust and power on rotor systems.


The conservation equations can be written by isolating  :

  \begin{enumerate}
    \item a large streamtube, englobing the full rotor section. This approach is commonly used in many situation,
    including in jet propulsion case.
    \item a small streamtube defined between $r$ and $r+dr$.
    It will be used section \ref{sec:thrust} to develop BEM theory.
  \end{enumerate}


%*****************************
\subsection{Large stream tube}
%******************************

\begin{figure}[h]
  \begin{center}
  \input{Figures/config_wind_turbine.pdf_tex}
  \caption{ Wind turbine rotor stream tube.}
  \label{fig:config_wind_turbine}
  \end{center}
\end{figure}



%**********************************************
\subsubsection{\sc Assumptions and mass flow rate }
%**********************************************
\label{sec:assumptions}

Now let us consider the first case  and notations of the figure  \ref{fig:config_wind_turbine}.
The conservation equations  are written twice between:

  \begin{itemize}
    \item the inlet section and the rotor plane
    \item the rotor plane and the outlet section 
  \end{itemize}


The main assumptions can be summerized as :

\begin{itemize}
\item steady incompressible 2D flow 
\item uniform velocity in a given section
\item viscosity and turbulence are not taken into account \textit{i.e.} laminar inviscid flow
\item no heat transfer 
\item the rotor is considered as a plane, with an infinite number of blades.
\item the inlet and outlet pressure are equal to atmospheric pressure (reference pressure)
\item thrust is uniformly distributed on the disk
\end{itemize}


The mass flow rate $q_m$ is given by :

\begin{equation}
  q_m = \rho S_0 V_0 = \rho S_1 V_1 = \rho S U_x
  \label{eq:qm}
\end{equation}

$U_x$ is the streamwise velocity across the rotor disk.

%**********************************************
\subsubsection{\sc Axial force}
%**********************************************
\label{sec:axial_force}

the axial force $T$ exerted by the rotor on the fluid is obtained from
the momentum equation projected on the streamwise direction : 
\begin{equation}
T = q_m (V_1 - V_0) =  \rho S U_x (V_1 - V_0) 
\label{eq:momentum_x}
\end{equation}


Let know consider the Bernoulli equation in the streamline arriving on the disk rotor from upstream,
and the same,  between the section just behind the rotor disk and the outlet section:

$$ \dsp 
p_0 +\rho  \frac{V_0^2}{2} = p^- + \rho  \frac{U_x^2}{2}, \qquad
p_1 +\rho  \frac{V_1^2}{2} = p^+ + \rho  \frac{U_x^2}{2}
$$

with $p_0=p_1=p_a$. The  axial force exerted by the rotor disk  on the fluid
is therefore :

\begin{equation} \dsp 
T = S \times (p^+-p^-) = \frac{S ~\rho}{2}~(V_1^2-V_0^2)= S ~\rho~(V_1-V_0)~ \frac{V_1+V_0}{2}
\label{eq:T_0}
\end{equation}

By comparing equation \eqref{eq:momentum_x} and \eqref{eq:T_0} it is found :
$$ \dsp
U_x = \frac{V_1+V_0}{2}  \qquad {\rm and} \qquad V_1 = 2 U_x -V_0
$$



An \textbf{axial induction factor} $a$ is usually defined to quantify the variation  of axial velocity across the 
rotor disk as :
\begin{equation} \dsp 
U_x = (1-a) V_0
\label{eq:def_a_0}
\end{equation}
$a$ is small and positive, the axial fluid velocity across the rotor  disk is smaller that the upstream velocity.
We also  get $$V_1 = V_0 (1-2 a)$$


Finally by substituting $V_1$ , it is found :

\begin{equation} \dsp 
T =  2~S~\rho~U_x~(U_x-V_0) = - T_{fluid\to rotor}
\end{equation}


The thrust coefficient is defined from $T$ as :

$$
\dsp
C_T = \frac{T_{fluid\to rotor}}{q S} = \frac{-T}{q S} = - 4 \epsilon (\epsilon -1) = 4 a (1- a)
$$


where the dynamic pressure is given by $q =\dsp \frac{1}{2} \rho V_0^2$, 
$\epsilon$ is  the nondimensionnal velocity $U_x/V_0$ and $\epsilon = 1 - a$

The maximun thrust coefficient is obtained for $a=1/2$ and $C_{T_{\max}} = 1$.

We can also express $a$ from the $C_T$:
$$ \dsp
a = \frac{1\pm\sqrt{1-C_T}}{2}
$$

For a given value of the $C_T$ we can be found two values of $a$, the first one in $[0, 1/2]$
and the second one in $[1/2, 1]$, as it shown on figure \ref{fig:C_T}.

\begin{figure}[h]
  \begin{center}
  \input{Figures/C_T.pdf_tex}
  \caption{ Wind turbine : $C_T$ and $C_P$ function of $a$.}
  \label{fig:C_T}
  \end{center}
\end{figure}

%***************************************************
\subsubsection{\sc Available power and efficiency}
%**************************************************

The power extracted (available power) from the actuator disk is given by:
$$
P_a = T_{fluid\to rotor} \times U_x  = -2~S~\rho~U_x^2~(U_x-V_0) 
$$

The power coefficient is a non dimensional coefficient defined as:

\begin{equation} \dsp 
  C_P = \frac{P_a}{q S V_0} = - 4 \epsilon^2 (\epsilon -1) = 4 a (1- a)^2
  \label{eq:Pa_0}
\end{equation}

It can easily be shown that $C_P$ is maximun for $a=1/3$, and $C_{P_{\max}} = 16/27\approx 0.593$ (see fig. \ref{fig:C_T}).

This coefficient can be interpreted as an efficiency, where $q S V_0$ should be the maximun
power provided by the wind. It is called the Betz limit. 


%********************************************
\subsubsection{\sc Torque }
%********************************************

The torque $Q$ exerted by the fluid on the rotor disk is calculated from the available power:

$$ 
Q \times \Omega = P_a = T_{fluid\to rotor} \times U_x \qquad  \Longrightarrow \qquad
Q =  \frac{T_{fluid\to rotor} \times U_x}{\Omega} = q S \frac{V_0}{\Omega} C_P
$$

Considering $R$ as the radius of the rotor disk ($S = \pi R^2$) then the non dimensional torque coefficient can
be written as:

$$\dsp
C_Q = \frac{Q}{q S R} = \frac{V_0}{R \Omega} C_P = \frac{C_P}{\lambda_r}
$$

In the following we will define the reduced velocity $\lambda_R = \dsp \frac{R~\Omega}{V_0}$.

Usually the reduced velocity is defined by the \textbf{advanced ratio parameter} $J$:

$$\dsp J = \frac{V_0}{N D} = \frac{V_0}{2 R \frac{\pi}{30} \Omega} = \dfrac{\pi}{15 \lambda_r} $$
where $N$ is the rotation velocity in rpm (rounds per minute).

and 


$$\dsp
C_Q = \frac{15 J}{\pi} C_P
$$

\Note{
One can found different formulas if we use $N$ , $D=2 \times R$ and $\rho V_0^2$ instead of $\Omega$, $R$ and $q$
as reference quantities. In that case the coefficients can be named $k_T$, $k_Q$ and $k_P$ instead of
$C_T$, $C_Q$ and $C_P$.}

 


%***********************************************
\subsection{Momentum theory with wake rotation}
%***********************************************



%**********************************
\subsubsection{\sc New assumption}
%**********************************

\begin{figure}[h]
  \begin{center}
  \input{Figures/disk_rotor_scheme.pdf_tex}
  \caption{Scheme of stream tube englobing the disk rotor (left), 
  in blue a streamline when rotational wake is added. On the right,
  a cut in a ($x, r$) plane, with useful notations.}
  \label{fig:disk_rotor_scheme}
  \end{center}
\end{figure}

The previous theory assumed that the rotor disk does not generate any rotation in the wake, which is not exactly true.
Therefore the momentum theory has been modified to allow flow rotation in the wake, assuming:

\begin{itemize}
  \item the upstream flow is not affected by the rotation
  \item immediatly behind the rotor section a tangential flow is generated and related to the
    downstream wake by an angular induction factor :
    \begin{equation} \dsp
      a' = \frac{\omega}{\Omega}
    \label{eq:ap0}
    \end{equation} 
    \todo{Verify if not $a' = \frac{\omega}{2\Omega}$, may be confusion between propeller and turbine}

    where $\Omega$ is the angular velocity of the rotor, $2 \times \omega$ 
    is the angular velocity imparted by the wake, and related to 
    wake vorticity with the additional assumption $\omega << \Omega$.
\end{itemize}

On figure \ref{fig:disk_rotor_scheme} is illustrated the effect of rotation assumption on a streamline (in blue).

%*******************************
\subsubsection{ \sc Thrust}
%*******************************
\label{sec:thrust}

The thrust calculation is also refered as the  momentum theory.


Let's considerer the rotor disk as infinitly thin and a thin stream tube of thickness $dr$ which provides
an axial force $dT$ on a small part of the blades, as depicted in figure \ref{fig:disk_rotor_scheme}.

From the momentum equation  (similarly to section \label{sec:axial_force})  it is found that :

\begin{equation} \dsp 
dT = dq_m (V_1(r) - V_0(r)) = \rho~dS~U_x~(V_1 -V_0 )
\label{eq:1}
\end{equation}

$dq_m$ is the elementary mass flow rate accross the stream tube, $dS = 2 \pi r dr$ is the section of the tube
at the location of the rotor disk, and $V_0$ and $V_1$ are the inlet and outlet tube velocities.
$U_x$ remains the fluid velocity across the rotor disk 

Let know consider the Bernoulli equation in the streamline arriving on the disk rotor,
and another one which starts from this section:

$$ \dsp 
p_0 +\rho  \frac{V_0^2}{2} = p^- + \rho  \frac{U_x^2}{2}, \qquad
p_1 +\rho  \frac{V_1^2}{2} = p^+ + \rho  \frac{U_x^2}{2}
$$

with $p_0=p_1$. The elementary axial force exerted by the disk rotor on the fluid
is therefore :

\begin{equation} \dsp 
dT = dS \times (p^+-p^-) = \frac{dS ~\rho}{2}~(V_1^2-V_0^2)= dS ~\rho~(V_1-V_0)~ \frac{V_1+V_0}{2}
\label{eq:2}
\end{equation}

By comparing equation \eqref{eq:1} and \eqref{eq:2} it is found :
$$ \dsp
\dsp U_x = \frac{V_1+V_0}{2}  \qquad {\rm and} \qquad V_1 = 2 U_x -V_0
$$

The relation is the same as the one found in the previous section, but here the velocities depends on the
radial coordinate $r$.

Finally by substituting $V_1$ , it is found :

\begin{equation} \dsp 
dT =  2~dS~\rho~U_x~(U_x-V_0) = 4~\pi~r~\rho~U_x~(U_x-V_0)~dr
\label{eq:3}
\end{equation}


The \textbf{axial induction factor} $a$ can always be defined    as :

\begin{equation} \dsp 
U_x(r) = (1-a(r)) V_0(r)
  \label{eq:def_ap}
\end{equation}

and the elementary axial force becomes
$\dsp  dT = - 4~\pi~r~\rho~V_0^2~a~(1-a)~dr$

and finally

\begin{equation}
dT_{fluid\to rotor}(r) = 4~\pi~r~\rho~V_0^2~a~(1-a)~dr
\label{eq:Tf-r}
\end{equation}
 



%********************************************
\subsubsection{\sc Tangential velocity and torque}
%********************************************
\label{secsec:ut}

The theory to determine the torque is also known as Tangential Momentum Theory.

Before the stream flow goes trough the rotor disk, no rotation is observed.
Across the rotor disk, in the rotation based at angular  velocity $\Omega$,
the fluid rotates at a different angular velocity $\omega$, opposite at the rotation direction of the rotor,
leading to a tangential perturbation velocity $u_t$, small, but proportional to $r~\Omega$.
Considering a  stream tube (cf. fig. \ref{fig:disk_rotor_scheme}),
tangential velocity perturbation is defined by   
a \textbf{tangential induction factor} $a'$

\begin{equation}
u_t = r ~\omega = a'(r)~r~\Omega 
\label{eq:def_aprime}
\end{equation}

In the outlet section of the stream tube, Glauert assumes that the angular velocity is $2\omega$,
and therefore the perturbation of tangential velocity is approximately twice 
the perturbation velocity in the disk rotor section. 

The variation of torque $dQ_{fluid\to rotor}=dQ$ due to the variation of tangential velocity is given by the second
Newton's law for the moment, written between the inlet and the outlet of the stream tube :

$$
dQ = \rho~U_x~r~u_{t_{outlet}}~dS= dq_m~r~u_{t_{outlet}} =  \rho~U_x~2\pi~r~dr~r~2~u_t
$$

Finaly by introducing the tangential and axial induction factor it yields :

\begin{equation}
dQ = 4~\pi~\rho~r^3~V_0(r)~\Omega~a'(1-a)~dr
\label{eq:def_dM}
\end{equation}



%********************************************
\subsubsection{\sc Power and propulsive efficiency}
%********************************************


The efficiency is the ratio between the output power to the input power, on the small stream tube of thickness $dr$ :

$$ \dsp
\eta =  \frac{\int \Omega~dQ}{\int V_0~dT_{fluid\to rotor}}  
$$
 


%****************************************
\section{Blade Element Momentum Theory}
%****************************************

%*************************
\subsection{Methodology}
%*************************

Different steps are to be developped:

\begin{enumerate}
\item blade geometry and discretisation
\item characteristic of the upstream flow
\item velocity decomposition
\item forces and torque
\item power and efficiency
\end{enumerate}

The velocity and forces section are related to aerodynamics.

When all these points are treated it remains to propose an algorithm to solve
the problem numerically and to discuss about options.

Results and conclusion will end the chapter.



%*****************************************
\subsection{Geometry}
%*****************************************
\label{sec:wt_gmtry}

\begin{figure}
  \begin{center}
  \input{Figures/rotor_and_propeller.pdf_tex}
  \caption{Rotor disk  of a wind turbine and propeller of an aircraft.}
  \label{fig:rotor_and_propeller}
  \end{center}
  \end{figure}

The description of the geometry is the same for a propeller and and wind turbine rotor.

The propeller and the rotor are build with $N_b$ blades distributed equally spaced (optionally) in the rotor section 
(fig. \ref{fig:rotor_and_propeller}). 
Blade  geometry is defined in figure \ref{fig:blades}.

The blade is discretized into $n$ elements, along the radial direction, each element is located between $r_k$ and $r_{k+1}$
for $k \in \{ 0, n \}$ and the zone of interest is for $r \in \{r_0, R\}$.

The airfoil section changes along the blade span so the blade is divided into $n_p$ panels, each panel is assumed to have the same airfoil section.

The blade is also defined by two laws:
\begin{itemize}
  \item the chord law $c(r)$
  \item the pitching angle law $\theta(r)$
\end{itemize}


The local polar base is $(\bs{e_x},\bs{e_r},\bs{e_t})$ (fig. \ref{fig:rotor_and_propeller}).

\begin{figure}
\begin{center}
\input{Figures/wind_turbine_blade.pdf_tex}
\caption{Schemes of wind rotor(bottom) and modern propeller (top) blades.}
\label{fig:blades}
\end{center}
\end{figure}

% \begin{figure}
% \begin{center}
% %****************************************
\section{Introduction}
%****************************************

The BEM theory has been first developped for propellers, used in boat and
aircrafts. Application to wind mills came some years later with Betz.

A wind turbine usually, does not moved, or at very low velocity, and the upstream wind is always
quite small, with a possible shear in the upstream wind.

At the opposite, propellers are used for aircraft, flying sometimes at more than Mach number equal to 0.5
and the upstream velocity profile is considered constant because of the small diameter of the propellers (few meters).

Another important difference is the shape of the stream tube englobing the rotor. Invertely to wind turbine, the 
tube section decreases with the streamwise direction leading to an acceletation of the flow.
To get this effect which is necessary to get a postive thrust on the aircraft, the pitch of the airfoil along
the blade is symmetrical with respect to the airfoil axis, comparing to a wind turbine.

All these differences are translated in the BEMT by some changes described in this chapter. 
Many details will not be given, the reader will be refer to the previous chapter.

%***********************************************************
\section{Non dimensional coefficients for propeller design}
%***********************************************************

The notation for propeller can be different from the those for wind turbine.
The main parameters for non dimensionnalisation are :
\begin{itemize}
\item $n$ : angular velocity in round/s
\item $V_0$ : upstream (aircraft) velocity
\item $D$ : propeller diameter
\item $\rho$ : air density
\item $q = \frac{1}{2} \rho V_0^2$
\end{itemize}

\Important{
$$ \begin{array}{llcl}
\textrm{Advance ratio}                &J &=& \dfrac{V_0}{nD}  \vs{0.3}\\
\textrm{Thrust}                       &T &=& \rho~D^4~n^2~k_T   \vs{0.3}\\
\textrm{Torque}                       &Q &=& \rho~D^5~n^2~k_Q   \vs{0.3}\\
\textrm{useful power}                 &P_u &=& \rho~D^5~n^3~k_{P_u} = T \times V_0 \qquad \Longrightarrow \qquad k_{P_u} = J \times k_T \vs{0.3}\\
\textrm{required power}               &P   &=& \rho~D^5~n^3~k_{P}   = Q \times \Omega \qquad \Longrightarrow \qquad k_{P} = 2\pi \times k_Q \vs{0.3}\\
\textrm{Propulsion efficiency} \hs{1} &\eta &=& \dfrac{P_u}{P} = \dfrac{T \times V_0}{Q \times \Omega} = \dfrac{k_T}{k_Q} \times \dfrac{J}{2\pi}    
\end{array}$$
}

\begin{itemize}
  \item The required power for the propeller is usually refered as airscrew power. 
  \item The useful power $P_u$ is the available power for the aircraft refers as $P_a$.
  \item we can also define the non dimensional coefficient as : 
  $$ \begin{array}{lcl}
    C_T &=& \dfrac{T}{q~S}  \vs{0.3}\\
    C_Q &=& \dfrac{Q}{q~S~D} \vs{0.3}\\
    C_P &=& \dfrac{Q}{q~S~V_0}
  \end{array}$$ 
\end{itemize}
%****************************************
\section{Rotor disk theory}
%****************************************

The two approaches developped in the previous chapter  will also be reported here, 
under a simplified form.



%*****************************
\subsection{Large stream tube}
%******************************

\begin{figure}[htb]
  \begin{center}
  \input{Figures/config_propeller.pdf_tex}
  \caption{Propeller rotor stream tube.}
  \label{fig:config_propeller}
  \end{center}
\end{figure}


%**********************************************
\subsubsection{\sc Assumptions and mass flow rate }
%**********************************************

see section \ref{sec:assumptions}



%**********************************************
\subsubsection{\sc Axial force}
%**********************************************
Developpement are identical to section \ref{sec:axial_force}.

However, because of the tube stream geometry, the definition of the
\textbf{axial induction factor} $a$ is slightly modified :

\begin{equation} \dsp 
U_x = (1+a) V_0
\label{eq:def_a_0_prop}
\end{equation}

$a$ is small and positive, the axial fluid velocity across the rotor  disk is greater than the upstream velocity.
We also  get $$V_1 = V_0 (1+2 a)$$

Comparing to wing turbine we just have to change the sign of $a$ in the formulas.

Finally by substituting $V_1$ , it is found :

\begin{equation} \dsp 
T =  2~S~\rho~U_x~(U_x-V_0) = - T_{fluid\to rotor}
\end{equation}


The thrust coefficient is defined from $T$ as :

$$
\dsp
C_T = \frac{T_{fluid\to rotor}}{q S} = \frac{-T}{q S} = - 4 \epsilon (\epsilon -1) = - 4 a (1+ a) < 0
$$



where the dynamic pressure is given by $q =\dsp \frac{1}{2} \rho V_0^2$, 
$\epsilon$ is  the nondimensionnal velocity $U_x/V_0$ and $\epsilon = 1 + a$



\Note{
\begin{itemize}
\item The axial force which is a thrust has changed of sign comparing to wind turbine.
\item For simplification, in the following we will forget the sign, and the axial force is assumed to be in $-x$ direction.
\end{itemize}}


So now we have $C_T = 4 a (1+ a)$.

We can   express $a$ from the $C_T$:
$$ \dsp
a = \frac{-1+\sqrt{1+C_T}}{2} = \dfrac{U_x-V_0}{V_0}
$$
Here we only have a single possible solution, since $a$ is a positive number.


%***************************************************
\subsubsection{\sc Available power and efficiency}
%**************************************************

\begin{itemize}
\item The power provided by the flow to the propeller is given by
$$
P = T \times U_x
$$
\item The power available for the aircraft propulsion is :
$$ 
P_a = T \times V_0 
$$
\item the induced power is given from the induced velocity $U_x - V_0$:
$$ P_i = T \times (U_x - V_0) $$
\end{itemize}


The \textbf{ideal efficiency} is defined by
$$\dsp
\eta_i = \frac{P_a}{P} = \frac{T \times V_0}{T \times U_x} = \frac{V_0}{(1+a) V_0} = \frac{1}{1+a}
= \frac{2}{1+\sqrt{1+C_T}}
$$


The power coefficient is a non dimensional coefficient defined as:

\begin{equation} \dsp 
  C_P = \frac{P}{q S V_0} = 4 a (1+ a)^2 = C_T~(1+a) = (1+\sqrt{1+C_T}) \times \frac{C_T}{2}
  \label{eq:Pa_0_prop}
\end{equation}




%********************************************
\subsubsection{\sc Torque }
%********************************************

The torque $Q$ provided by the engine is calculated from the  power provided by the flow $P$ :

$$ 
Q \times \Omega = P = T \times U_x \qquad  \Longrightarrow \qquad
Q =  \frac{P}{\Omega} = q S \frac{V_0}{\Omega} C_P
$$

Considering $D$ as the diameter of the rotor disk ($S = \pi D^2/4$) then the non dimensional torque coefficient can
be written as:

$$\dsp
C_Q = \frac{Q}{q S D} = \frac{V_0}{D \Omega} C_P = \frac{C_T}{2~\lambda_R}
$$

where $\lambda_R$  the reduced velocity $\lambda_R = \dsp \frac{R~\Omega}{V_0}$.

Usually the reduced velocity is related to  \textbf{advance ratio parameter} $J$:
$$
J \times \lambda_R = \pi
$$






%***********************************************
\subsection{Momentum theory with wake rotation}
%***********************************************

The theory does not differ from the one given in the previous chapter.
The difference is the changed of signe in front of $a$.
Therefore we get for the elementary thrust \eqref{eq:Tf-r} and elementary torque the equations:


\begin{eqnarray}
  dT(r) &=& 4~\pi~r~\rho~V_0^2~a~(1+a)~dr \label{eq:Tf-r2}  \vs{0.5} \\
  dQ(r) &=& 4~\pi~\rho~r^3~V_0~\Omega~a'(1+a)~dr \label{eq:def_dM2}
\end{eqnarray}

  


%****************************************
\section{Blade Element Momentum Theory}
%****************************************


%*****************************************
\subsection{Geometry}
%*****************************************

%**************************************
\subsubsection{\sc Pitch of propeller}
%**************************************

\todo{See \cite{VonMises1959} pages 289-295, for the geometry of a propeller and the
calculation of the pitch angle}

%*******************************
\subsubsection{\sc Discretisation}
%******************************

\begin{figure}
  \begin{center}
  %****************************************
\section{Introduction}
%****************************************

The BEM theory has been first developped for propellers, used in boat and
aircrafts. Application to wind mills came some years later with Betz.

A wind turbine usually, does not moved, or at very low velocity, and the upstream wind is always
quite small, with a possible shear in the upstream wind.

At the opposite, propellers are used for aircraft, flying sometimes at more than Mach number equal to 0.5
and the upstream velocity profile is considered constant because of the small diameter of the propellers (few meters).

Another important difference is the shape of the stream tube englobing the rotor. Invertely to wind turbine, the 
tube section decreases with the streamwise direction leading to an acceletation of the flow.
To get this effect which is necessary to get a postive thrust on the aircraft, the pitch of the airfoil along
the blade is symmetrical with respect to the airfoil axis, comparing to a wind turbine.

All these differences are translated in the BEMT by some changes described in this chapter. 
Many details will not be given, the reader will be refer to the previous chapter.

%***********************************************************
\section{Non dimensional coefficients for propeller design}
%***********************************************************

The notation for propeller can be different from the those for wind turbine.
The main parameters for non dimensionnalisation are :
\begin{itemize}
\item $n$ : angular velocity in round/s
\item $V_0$ : upstream (aircraft) velocity
\item $D$ : propeller diameter
\item $\rho$ : air density
\item $q = \frac{1}{2} \rho V_0^2$
\end{itemize}

\Important{
$$ \begin{array}{llcl}
\textrm{Advance ratio}                &J &=& \dfrac{V_0}{nD}  \vs{0.3}\\
\textrm{Thrust}                       &T &=& \rho~D^4~n^2~k_T   \vs{0.3}\\
\textrm{Torque}                       &Q &=& \rho~D^5~n^2~k_Q   \vs{0.3}\\
\textrm{useful power}                 &P_u &=& \rho~D^5~n^3~k_{P_u} = T \times V_0 \qquad \Longrightarrow \qquad k_{P_u} = J \times k_T \vs{0.3}\\
\textrm{required power}               &P   &=& \rho~D^5~n^3~k_{P}   = Q \times \Omega \qquad \Longrightarrow \qquad k_{P} = 2\pi \times k_Q \vs{0.3}\\
\textrm{Propulsion efficiency} \hs{1} &\eta &=& \dfrac{P_u}{P} = \dfrac{T \times V_0}{Q \times \Omega} = \dfrac{k_T}{k_Q} \times \dfrac{J}{2\pi}    
\end{array}$$
}

\begin{itemize}
  \item The required power for the propeller is usually refered as airscrew power. 
  \item The useful power $P_u$ is the available power for the aircraft refers as $P_a$.
  \item we can also define the non dimensional coefficient as : 
  $$ \begin{array}{lcl}
    C_T &=& \dfrac{T}{q~S}  \vs{0.3}\\
    C_Q &=& \dfrac{Q}{q~S~D} \vs{0.3}\\
    C_P &=& \dfrac{Q}{q~S~V_0}
  \end{array}$$ 
\end{itemize}
%****************************************
\section{Rotor disk theory}
%****************************************

The two approaches developped in the previous chapter  will also be reported here, 
under a simplified form.



%*****************************
\subsection{Large stream tube}
%******************************

\begin{figure}[htb]
  \begin{center}
  \input{Figures/config_propeller.pdf_tex}
  \caption{Propeller rotor stream tube.}
  \label{fig:config_propeller}
  \end{center}
\end{figure}


%**********************************************
\subsubsection{\sc Assumptions and mass flow rate }
%**********************************************

see section \ref{sec:assumptions}



%**********************************************
\subsubsection{\sc Axial force}
%**********************************************
Developpement are identical to section \ref{sec:axial_force}.

However, because of the tube stream geometry, the definition of the
\textbf{axial induction factor} $a$ is slightly modified :

\begin{equation} \dsp 
U_x = (1+a) V_0
\label{eq:def_a_0_prop}
\end{equation}

$a$ is small and positive, the axial fluid velocity across the rotor  disk is greater than the upstream velocity.
We also  get $$V_1 = V_0 (1+2 a)$$

Comparing to wing turbine we just have to change the sign of $a$ in the formulas.

Finally by substituting $V_1$ , it is found :

\begin{equation} \dsp 
T =  2~S~\rho~U_x~(U_x-V_0) = - T_{fluid\to rotor}
\end{equation}


The thrust coefficient is defined from $T$ as :

$$
\dsp
C_T = \frac{T_{fluid\to rotor}}{q S} = \frac{-T}{q S} = - 4 \epsilon (\epsilon -1) = - 4 a (1+ a) < 0
$$



where the dynamic pressure is given by $q =\dsp \frac{1}{2} \rho V_0^2$, 
$\epsilon$ is  the nondimensionnal velocity $U_x/V_0$ and $\epsilon = 1 + a$



\Note{
\begin{itemize}
\item The axial force which is a thrust has changed of sign comparing to wind turbine.
\item For simplification, in the following we will forget the sign, and the axial force is assumed to be in $-x$ direction.
\end{itemize}}


So now we have $C_T = 4 a (1+ a)$.

We can   express $a$ from the $C_T$:
$$ \dsp
a = \frac{-1+\sqrt{1+C_T}}{2} = \dfrac{U_x-V_0}{V_0}
$$
Here we only have a single possible solution, since $a$ is a positive number.


%***************************************************
\subsubsection{\sc Available power and efficiency}
%**************************************************

\begin{itemize}
\item The power provided by the flow to the propeller is given by
$$
P = T \times U_x
$$
\item The power available for the aircraft propulsion is :
$$ 
P_a = T \times V_0 
$$
\item the induced power is given from the induced velocity $U_x - V_0$:
$$ P_i = T \times (U_x - V_0) $$
\end{itemize}


The \textbf{ideal efficiency} is defined by
$$\dsp
\eta_i = \frac{P_a}{P} = \frac{T \times V_0}{T \times U_x} = \frac{V_0}{(1+a) V_0} = \frac{1}{1+a}
= \frac{2}{1+\sqrt{1+C_T}}
$$


The power coefficient is a non dimensional coefficient defined as:

\begin{equation} \dsp 
  C_P = \frac{P}{q S V_0} = 4 a (1+ a)^2 = C_T~(1+a) = (1+\sqrt{1+C_T}) \times \frac{C_T}{2}
  \label{eq:Pa_0_prop}
\end{equation}




%********************************************
\subsubsection{\sc Torque }
%********************************************

The torque $Q$ provided by the engine is calculated from the  power provided by the flow $P$ :

$$ 
Q \times \Omega = P = T \times U_x \qquad  \Longrightarrow \qquad
Q =  \frac{P}{\Omega} = q S \frac{V_0}{\Omega} C_P
$$

Considering $D$ as the diameter of the rotor disk ($S = \pi D^2/4$) then the non dimensional torque coefficient can
be written as:

$$\dsp
C_Q = \frac{Q}{q S D} = \frac{V_0}{D \Omega} C_P = \frac{C_T}{2~\lambda_R}
$$

where $\lambda_R$  the reduced velocity $\lambda_R = \dsp \frac{R~\Omega}{V_0}$.

Usually the reduced velocity is related to  \textbf{advance ratio parameter} $J$:
$$
J \times \lambda_R = \pi
$$






%***********************************************
\subsection{Momentum theory with wake rotation}
%***********************************************

The theory does not differ from the one given in the previous chapter.
The difference is the changed of signe in front of $a$.
Therefore we get for the elementary thrust \eqref{eq:Tf-r} and elementary torque the equations:


\begin{eqnarray}
  dT(r) &=& 4~\pi~r~\rho~V_0^2~a~(1+a)~dr \label{eq:Tf-r2}  \vs{0.5} \\
  dQ(r) &=& 4~\pi~\rho~r^3~V_0~\Omega~a'(1+a)~dr \label{eq:def_dM2}
\end{eqnarray}

  


%****************************************
\section{Blade Element Momentum Theory}
%****************************************


%*****************************************
\subsection{Geometry}
%*****************************************

%**************************************
\subsubsection{\sc Pitch of propeller}
%**************************************

\todo{See \cite{VonMises1959} pages 289-295, for the geometry of a propeller and the
calculation of the pitch angle}

%*******************************
\subsubsection{\sc Discretisation}
%******************************

\begin{figure}
  \begin{center}
  %****************************************
\section{Introduction}
%****************************************

The BEM theory has been first developped for propellers, used in boat and
aircrafts. Application to wind mills came some years later with Betz.

A wind turbine usually, does not moved, or at very low velocity, and the upstream wind is always
quite small, with a possible shear in the upstream wind.

At the opposite, propellers are used for aircraft, flying sometimes at more than Mach number equal to 0.5
and the upstream velocity profile is considered constant because of the small diameter of the propellers (few meters).

Another important difference is the shape of the stream tube englobing the rotor. Invertely to wind turbine, the 
tube section decreases with the streamwise direction leading to an acceletation of the flow.
To get this effect which is necessary to get a postive thrust on the aircraft, the pitch of the airfoil along
the blade is symmetrical with respect to the airfoil axis, comparing to a wind turbine.

All these differences are translated in the BEMT by some changes described in this chapter. 
Many details will not be given, the reader will be refer to the previous chapter.

%***********************************************************
\section{Non dimensional coefficients for propeller design}
%***********************************************************

The notation for propeller can be different from the those for wind turbine.
The main parameters for non dimensionnalisation are :
\begin{itemize}
\item $n$ : angular velocity in round/s
\item $V_0$ : upstream (aircraft) velocity
\item $D$ : propeller diameter
\item $\rho$ : air density
\item $q = \frac{1}{2} \rho V_0^2$
\end{itemize}

\Important{
$$ \begin{array}{llcl}
\textrm{Advance ratio}                &J &=& \dfrac{V_0}{nD}  \vs{0.3}\\
\textrm{Thrust}                       &T &=& \rho~D^4~n^2~k_T   \vs{0.3}\\
\textrm{Torque}                       &Q &=& \rho~D^5~n^2~k_Q   \vs{0.3}\\
\textrm{useful power}                 &P_u &=& \rho~D^5~n^3~k_{P_u} = T \times V_0 \qquad \Longrightarrow \qquad k_{P_u} = J \times k_T \vs{0.3}\\
\textrm{required power}               &P   &=& \rho~D^5~n^3~k_{P}   = Q \times \Omega \qquad \Longrightarrow \qquad k_{P} = 2\pi \times k_Q \vs{0.3}\\
\textrm{Propulsion efficiency} \hs{1} &\eta &=& \dfrac{P_u}{P} = \dfrac{T \times V_0}{Q \times \Omega} = \dfrac{k_T}{k_Q} \times \dfrac{J}{2\pi}    
\end{array}$$
}

\begin{itemize}
  \item The required power for the propeller is usually refered as airscrew power. 
  \item The useful power $P_u$ is the available power for the aircraft refers as $P_a$.
  \item we can also define the non dimensional coefficient as : 
  $$ \begin{array}{lcl}
    C_T &=& \dfrac{T}{q~S}  \vs{0.3}\\
    C_Q &=& \dfrac{Q}{q~S~D} \vs{0.3}\\
    C_P &=& \dfrac{Q}{q~S~V_0}
  \end{array}$$ 
\end{itemize}
%****************************************
\section{Rotor disk theory}
%****************************************

The two approaches developped in the previous chapter  will also be reported here, 
under a simplified form.



%*****************************
\subsection{Large stream tube}
%******************************

\begin{figure}[htb]
  \begin{center}
  \input{Figures/config_propeller.pdf_tex}
  \caption{Propeller rotor stream tube.}
  \label{fig:config_propeller}
  \end{center}
\end{figure}


%**********************************************
\subsubsection{\sc Assumptions and mass flow rate }
%**********************************************

see section \ref{sec:assumptions}



%**********************************************
\subsubsection{\sc Axial force}
%**********************************************
Developpement are identical to section \ref{sec:axial_force}.

However, because of the tube stream geometry, the definition of the
\textbf{axial induction factor} $a$ is slightly modified :

\begin{equation} \dsp 
U_x = (1+a) V_0
\label{eq:def_a_0_prop}
\end{equation}

$a$ is small and positive, the axial fluid velocity across the rotor  disk is greater than the upstream velocity.
We also  get $$V_1 = V_0 (1+2 a)$$

Comparing to wing turbine we just have to change the sign of $a$ in the formulas.

Finally by substituting $V_1$ , it is found :

\begin{equation} \dsp 
T =  2~S~\rho~U_x~(U_x-V_0) = - T_{fluid\to rotor}
\end{equation}


The thrust coefficient is defined from $T$ as :

$$
\dsp
C_T = \frac{T_{fluid\to rotor}}{q S} = \frac{-T}{q S} = - 4 \epsilon (\epsilon -1) = - 4 a (1+ a) < 0
$$



where the dynamic pressure is given by $q =\dsp \frac{1}{2} \rho V_0^2$, 
$\epsilon$ is  the nondimensionnal velocity $U_x/V_0$ and $\epsilon = 1 + a$



\Note{
\begin{itemize}
\item The axial force which is a thrust has changed of sign comparing to wind turbine.
\item For simplification, in the following we will forget the sign, and the axial force is assumed to be in $-x$ direction.
\end{itemize}}


So now we have $C_T = 4 a (1+ a)$.

We can   express $a$ from the $C_T$:
$$ \dsp
a = \frac{-1+\sqrt{1+C_T}}{2} = \dfrac{U_x-V_0}{V_0}
$$
Here we only have a single possible solution, since $a$ is a positive number.


%***************************************************
\subsubsection{\sc Available power and efficiency}
%**************************************************

\begin{itemize}
\item The power provided by the flow to the propeller is given by
$$
P = T \times U_x
$$
\item The power available for the aircraft propulsion is :
$$ 
P_a = T \times V_0 
$$
\item the induced power is given from the induced velocity $U_x - V_0$:
$$ P_i = T \times (U_x - V_0) $$
\end{itemize}


The \textbf{ideal efficiency} is defined by
$$\dsp
\eta_i = \frac{P_a}{P} = \frac{T \times V_0}{T \times U_x} = \frac{V_0}{(1+a) V_0} = \frac{1}{1+a}
= \frac{2}{1+\sqrt{1+C_T}}
$$


The power coefficient is a non dimensional coefficient defined as:

\begin{equation} \dsp 
  C_P = \frac{P}{q S V_0} = 4 a (1+ a)^2 = C_T~(1+a) = (1+\sqrt{1+C_T}) \times \frac{C_T}{2}
  \label{eq:Pa_0_prop}
\end{equation}




%********************************************
\subsubsection{\sc Torque }
%********************************************

The torque $Q$ provided by the engine is calculated from the  power provided by the flow $P$ :

$$ 
Q \times \Omega = P = T \times U_x \qquad  \Longrightarrow \qquad
Q =  \frac{P}{\Omega} = q S \frac{V_0}{\Omega} C_P
$$

Considering $D$ as the diameter of the rotor disk ($S = \pi D^2/4$) then the non dimensional torque coefficient can
be written as:

$$\dsp
C_Q = \frac{Q}{q S D} = \frac{V_0}{D \Omega} C_P = \frac{C_T}{2~\lambda_R}
$$

where $\lambda_R$  the reduced velocity $\lambda_R = \dsp \frac{R~\Omega}{V_0}$.

Usually the reduced velocity is related to  \textbf{advance ratio parameter} $J$:
$$
J \times \lambda_R = \pi
$$






%***********************************************
\subsection{Momentum theory with wake rotation}
%***********************************************

The theory does not differ from the one given in the previous chapter.
The difference is the changed of signe in front of $a$.
Therefore we get for the elementary thrust \eqref{eq:Tf-r} and elementary torque the equations:


\begin{eqnarray}
  dT(r) &=& 4~\pi~r~\rho~V_0^2~a~(1+a)~dr \label{eq:Tf-r2}  \vs{0.5} \\
  dQ(r) &=& 4~\pi~\rho~r^3~V_0~\Omega~a'(1+a)~dr \label{eq:def_dM2}
\end{eqnarray}

  


%****************************************
\section{Blade Element Momentum Theory}
%****************************************


%*****************************************
\subsection{Geometry}
%*****************************************

%**************************************
\subsubsection{\sc Pitch of propeller}
%**************************************

\todo{See \cite{VonMises1959} pages 289-295, for the geometry of a propeller and the
calculation of the pitch angle}

%*******************************
\subsubsection{\sc Discretisation}
%******************************

\begin{figure}
  \begin{center}
  \input{Figures/propeller.pdf_tex}
  \caption{Rotor disk  of a propeller of the Airbus A400 M.}
  \label{fig:propeller_alone}
  \end{center}
\end{figure}

A example of modern aircraft propeller is shown in figure \ref{fig:propeller_alone}.

For additional details read the section \ref{sec:wt_gmtry}


%******************************
\subsection{Upstream conditions}
%******************************

As mentioned before, the upstream velocity is assumed constant: $V(r) = V_0$ in any equation.


%***********************
\subsection{Velocities}
%***********************

\begin{figure}[htb]
  \begin{center}
  \input{Figures/propeller_side_notations.pdf_tex}
  \caption{Scheme of a modern propeller, side view (left) and airfoil section with velocity and force decompositions (right)}
  \label{fig:propeller_side_notations}
  \end{center}
  \end{figure}
  
  To develop the BEM theory we have to refer to the figure \ref{fig:propeller_side_notations} which provides
  velocity and force decompositions.
  
  \Important{
  The main difference between a rotor blade section and a propeller blade section is the direction of the airfoil with respect
  to the rotation plane. The orientation on a propeller is choosen to produce positive thrust.
  }


  This orientation naturally modifies slightly the velocity and force directions and amplitudes as explained below.


  Let's consider a given airfoil section located at a distance $r$ from the axis of rotation (fig. \ref{fig:propeller_alone}). 
  The BEMT consists on analyzing the decomposition of the \textbf{apparent velocity} vector refered as $\bs{W}$ 
  which is the air velocity seen by the airfoil section and
  which allow to determine an angle of attack $\alpha$.
  
  The apparent velocity is the sum of :
  \begin{itemize}
  \item the constant upstream  velocity $\bs{V}$
  
  \item the tangential velocity related to the disk rotation $\bs{V_t} = - r~\Omega~\bs{e_t}$,
  where $\bs{e_t} = \bs{e_x} \times \bs{e_r}$. 
  
  \item the induced velocity $\bs{w}$ whose direction is defined by the induced angle of attack $\alpha_i$. 
  It is due to the tip vortex of the blade. Its direction is assumed orthogonal to the airfoil section axis.   
  \end{itemize}
  
  
  In addition we can define the relative velocity $\bs{V_R} = \bs{V} + \bs{V_t}$ whose direction is the angle $\varphi$.

  It can be easily demonstrated that :
  \begin{equation}
  \tan \varphi =  \frac{V}{r~\Omega} = \dfrac{J}{\pi \tilde r}, \qquad \tilde r = \dfrac{2~r}{D} 
\label{eq:varphi_prop}  
\end{equation}
  
  And finally, the apparent velocity can be projected into the $\bs{e_x}$ and $\bs{e_t}$ directions as :
  
  \begin{equation}
  \bs{W} = U_x~\bs{e_x} + U_t~\bs{e_t}
  \label{eq:W_def2}
  \end{equation} 
  
  where :
  \begin{itemize}
    \item $U_x(r)$ is the axial velocity, air velocity crossing the propeller plane
    \item $U_t(r)$ is the tangential velocity, in the propeller plane. 
  \end{itemize}
  
  All these velocities are drawn on figure  \ref{fig:propeller_side_notations}.

  \Note{
  In addition the tangential induced velocity $u_t = V_t - U_t$ is drawn in blue. Because of the direction
  of the airfoil section this velocity is opposite comparing to the rotor case, and because of the direction
  of the induced velocity $w$ we can notice that $U_x > V$.}

%********************
\subsection{Angles}
%********************

Angles play an important role to define velocity and force directions.

Different angles are defined from the velocities and aerodynamics :
\begin{itemize}
  \item the induced angle of attack $\alpha_i$
  \item the relative velocity angle $\varphi$
  \item the local flow angle $\phi$ with 
  $$
  \phi = \alpha_i  + \varphi
  $$
  \item the pitching angle $\theta$, which comes from the geometry
  \item the real angle of attack $\alpha$
  \item the aerodynamic efficiency angle $\gamma$ with $\tan \gamma = C_d/C_\ell = 1/f$.
\end{itemize}

$f$ is here the aerodynamic efficiency. $C_d$ and $C_{\ell}$ are the usual lift and drag aerodynamic 
coefficients  used in the next section.


To determine forces exerted on the airfoil section we need to clearly defined the angle of attack  by :

\Important{
\begin{equation}
\alpha = \theta-\phi
\label{eq:alpha_prop}
\end{equation}
}

\gris{We can see that the angle of attack equation is the opposite for a wind turbine rotor}

%************************************
\subsection{Forces on a wing section}
%************************************

On figure \ref{fig:propeller_side_notations} are drawn the aerodynamics force on a wing section.

Let write the equations for a given section, and for a single blade :

\begin{itemize}
    \item the lift 
    $$
      \dsp \bs{dL} = q~dS~C_\ell~\bs{j_a}  
    $$
    \item the drag 
    $$
      \dsp \bs{dD} = q~dS~C_d ~\bs{i_a}
    $$
    \item the aerodynamic force $\bs{dF}$
    $$\dsp
    \bs{dF} = dL~\bs{j_a} + dD~\bs{i_a} = \bs{dT} + \bs{dF_t}
    $$
    \item the normal aerodynamic force :
    $$
    dT = \bs{dF}\cdot\bs{e_x} = q~dS~C_n
    $$
    \item the tangential aerodynamic force:
    $$
    dF_t = \bs{dF}\cdot\bs{e_y} = q~dS~C_t
    $$

\end{itemize}

with 
\begin{itemize}
  \item $q=\frac{1}{2} \rho W^2$, the local dynamic pressure
  % \item $S = c(r) \times b(r) $ where $b(r)$ is the span of the local element.
  \item $C_\ell, C_d$ are the aerodynamic lift and drag coefficients dependent on the airfoil.
  \item $C_n, C_t$ are the aerodynamic normal and tangential  coefficient dependent on the $C_{\ell}$ and $C_d$.
  \item $(\bs{i_a}, \bs{j_a})$ is the local aerodynamic base ($\bs{W} = W~\bs{i_a}$) (see fig.~\ref{fig:rotor_side_notations})
  \item $ dS = c(r)~ dr$ is the area of a blade element
\end{itemize}

In addition it it easy to show that :
\begin{eqnarray}
dT &=& \cos\phi ~ dL - \sin\phi~ dD   \label{eq:dT}\\
dF_t &=& \sin\phi ~ dL + \cos\phi~ dD \label{eq:dFt} 
\end{eqnarray}

From equations \ref{eq:dT} and \ref{eq:dFt}  it is found :
\begin{equation}
\begin{array}{lcl}
C_n &=& C_\ell \cos\phi - C_d \sin \phi \\
C_t &=& C_\ell \sin\phi + C_d \cos \phi
\end{array}
\label{eq:CnCt}
\end{equation}


%********************************************
\subsection{BEM theory algorithm}
%********************************************

To develop the Blade Element Momentum Theory, several steps are necessary and some reduced quantities are 
introduced (ever discussed in the disk rotor theory).

%******************************************************************
\subsubsection{\sc Step 1 : induced factor and apparent velocity}
%*****************************************************************
The BEMT use different reduced quantities and the most important are the induced flow factor $a$ and $a'$ defined as :

\begin{itemize}
\item the induced axial velocity factor $a$ is given by :
\Def{induced axial velocity factor $a$}{
  \begin{equation}
  U_x~/~V = 1 + a
  \label{eq:a_def_prop}
  \end{equation}
}

From the figure \ref{fig:propeller_side_notations} it can be seen that $U_x > V$, that means that $0 < a $ and $a$ should be small.


\item the induced radial velocity factor $a'$ is given by 
\Def{induced radial velocity factor $a'$}{  
\begin{equation}
    U_t  = (1-a')~V_t = (1-a') ~r~\Omega
  \label{eq:aprime_def_prop}
  \end{equation}
}



\item We can then find :
\begin{equation} \dsp
\tan \phi = \frac{U_x}{U_t} = \frac{1 + a}{1-a'} \times \frac{V}{r~\Omega}, \qquad U_x = W \sin\phi, \qquad U_t = W \cos\phi
\label{eq:tanphi_prop}
\end{equation}
\end{itemize}

It is now possible to express the modulus of the apparent velocity as a function of the induced factors :
\begin{equation}
\dsp W = \sqrt{V^2(1+a)^2 + r^2 \Omega^2 (1-a')^2} = V  \sqrt{(1+a)^2 + \left(\frac{r \Omega}{V}\right)^2 (1-a')^2}
\label{eq:W_prop}
\end{equation}


From the previous equation, we can defined a reduced velocity parameter  refers as \textbf{local tip speed ratio}
and ever seen in equation \eqref{eq:varphi_prop}:
  
\begin{equation} \dsp
  \lambda_r =  \frac{r \Omega}{V} = \frac{1}{\tan \varphi}  = \frac{\pi~\tilde r}{J}
  \label{eq:lambda_r_prop}
\end{equation}


Finally, by using the previous equations, we can show the result:

\begin{equation} \dsp
\dfrac{a'}{a} = \frac{C_t}{\lambda_r C_n}
  \label{eq:aprime_over_a}
\end{equation}



%******************************************************
\subsubsection{\sc Step 2 : equations for induction factors}
%******************************************************

The next step is to make the relationship between velocity, induced factors 
and aerodynamic coefficients.

\begin{itemize}

\item  \textbf{axial induction factor} equation.

First we can write the elementary axial force   for $N_b$ blades as :
$$ \dsp
dT = N_b~q~dS~C_n = \frac{1}{2} \rho W^2 C_n  N_b~c(r)~dr, \quad {\rm with}\qquad U_x = W \sin\phi 
$$

The  $dT$ is also given from the rotor disk theory by the
equation \eqref{eq:Tf-r2}  

Equating the both  equations leads to :
\begin{equation}
\frac{C_n~N_b~c}{8 \pi r ~\sin^2\phi} =  \frac{a}{1+a}  
\label{eq:tmp1_prop}
\end{equation}


Here can be introduced the
\textbf{solidity function}  $\sigma$ (sometimes refered as $s$): 

\Def{Solidity $\sigma$}{
\begin{equation}
  \sigma = \dfrac{N_b~c(r)}{2\pi~r} = \dfrac{N_b~c(r)}{\pi~D} ~\frac{1}{\tilde r}
\end{equation}
}
where $c(r)$ is the chord law and $N_b$ the number of blades.


For propellers, as the one shown on figure   \ref{fig:rotor_and_propeller} the solidity is much larger, but
decreases with increasing $r$.



Finally the equation \eqref{eq:tmp1_prop} can be simplified 
by introducing the solidity parameter as  

\begin{equation} \dsp \dfrac{a}{1+a} = \dfrac{\sigma C_n}{4 \sin^2\phi}, \qquad
  \dfrac{1}{a} =  \dfrac{4 \sin^2\phi}{\sigma C_n}  - 1
  \label{eq:a_prop}\end{equation} 

\item  \textbf{tangential induction factor} equation :


The elementary torque on $N_b$ blades,
 from figure \ref{fig:rotor_side_notations} is given by 
$$
dQ = N_b~r~dF_t = N_b~r~q~c(r)~dr~C_t = \frac{\rho}{2} W^2~N_b~r~c(r)~dr~C_t
$$

In this equation we rewrite 
$\dsp W^2 = W\times W  =  \frac{U_x~U_t}{\cos\phi~\sin\phi}$ with
$U_t = r~\Omega~(1-a')$

It is equal to the elementary momentum given by the rotor disk theory equation
\eqref{eq:def_dM2} written as :
$$
dQ =  4 \  \pi \ \rho \ V \ \Omega \  a' \ (1+a)  \ r^3~dr= 4~\pi~\rho~r^3~\Omega~a'~U_x~dr
$$


finally we find :


\begin{equation} \dsp 
  \dfrac{a'}{1-a'}  =   \dfrac{\sigma C_t}{4 \sin\phi \cos\phi} \qquad
  \Longrightarrow \qquad \dfrac{1}{a'} = \dfrac{4 \sin\phi \cos\phi}{\sigma C_t}+1 
\label{eq:ap_prop}\end{equation}  

\item Local flow angle $\phi$, from the equation \eqref{eq:tanphi} we get
$$\sin\phi = \dfrac{V}{W} (1+a), \qquad  \cos\phi = \dfrac{r \Omega}{W} (1-a')
$$
\begin{equation}\lambda_r \tan \phi = \dfrac{1+a}{1-a'} \label{eq:lam_prop}\end{equation}  

\gris{This equation has ever been introduced in equation \eqref{eq:tanphi_prop}.}
\end{itemize}

A single equation without the induction factors $a$ and $a'$ can be obtained from 
equations \eqref{eq:a_prop}, \eqref{eq:ap_prop} and  \eqref{eq:lam_prop}

\Important{
  \begin{equation} \dsp 
      F_{\rm Newton}(\phi, r) = 0 =  - 1 +   \frac{1}{\lambda_r(r)} \ \frac{2 \sin2\phi +\sigma(r) C_t(\phi)}{4 \sin^2\phi - \sigma(r) C_n(\phi)} 
  \label{eq:Phi_prop}
  \end{equation} 
}

%******************************************************
\subsubsection{\sc Step 3 : thrust and torque}
%******************************************************

\begin{itemize}
  \item the elementary thust is given from previous section as :
  $$ \dsp
   \der{T}{r} =  \dfrac{1}{2}~\rho~N_b~V^2~c(r)~\dfrac{(1+a)^2}{\sin^2\phi}~C_n~=~
   4 \ \pi \ \rho \ V^2 \ a \ (1+a) \ r
  $$
  
  and the total  thrust is :
  $$ \dsp
  T  = \int_{r_0}^{D/2} dT
  $$
  
  \item the elementary torque has ever been calculated  as well:
  $$\dsp
  \der{Q}{r} =  \dfrac{1}{2}~\rho~N_b~(r~\Omega)^2~r~c(r)~\dfrac{(1-a')^2}{\cos^2\phi}~C_t~=~
   4~\pi~\rho~r^3~V~\Omega~a'(1+a)
  $$
  and the total torque is 
  
  $$ \dsp
  Q  = \int_{r_0}^{D/2} dQ
  $$
  \end{itemize}

%******************************************************
\subsubsection{\sc Step 4 : power and efficiency}
%******************************************************

\begin{itemize}
  \item the input propeller power is :
  $$ \dsp
  P_{in}  = \int_{r_0}^{D/2} \Omega\times dQ = \Omega \times Q 
  $$
  
  
  \item the output thrust power is 
    
  $$ \dsp
  P_{in}   = \int_{r_0}^{D/2} V \times dT = V \times T
  $$
  
  \item the efficiency is therefore:
  $$\dsp
  \eta = \frac{P_{out}}{P_{in}} = \frac{V}{\Omega} \times \frac{T}{Q} = \frac{V^2}{\Omega^2} \times
  \frac{\dsp \int_{r_0}^{D/2}  a (1+a) r~dr }{\dsp \int_{r_0}^{D/2} a'(1+a)r^3~dr} =
  \frac{J^2}{\pi^2} \times
  \frac{\dsp \int_{\tilde r_0}^{1}  a (1+a) \tilde r~d\tilde r }{\dsp \int_{\tilde r_0}^{1} a'(1+a)\tilde r^3~d\tilde r}
  $$
  \end{itemize}
  


%*******************************************************
\subsection{Tip loss}
%*******************************************************

In the previous theory the tip loss existing because the blade is a wing of finite span has not been taken
into account. Prandtl has proposed a correction  to the BEMT theory by introducing a function $F(r)$ as follows :

\begin{itemize}
\item for elementary torque and thrust :
$$ \dsp
\der{Q}{r} = 4 \  \pi \ \rho \ V \ \Omega \  a' \ (1+a)  \ r^3~F(\phi,r)
$$

$$ \dsp
\der{T}{r} =    4 \ \pi \ \rho \ V^2 \ a \ (1+a) \ r~F(\phi,r)
$$

with
$$ \dsp
F(\phi,r) = \dfrac{2}{\pi} \arccos {\rm e}^{-f} , \qquad f =  \dfrac{N_b}{2} \ \dfrac{R-r}{r \sin\phi})
$$

\item the induction factor equations are modified  as :

\begin{equation}
\begin{array}{lcl}
\dsp \dfrac{1}{a} &=& \dsp \dfrac{4 F \sin^2\phi}{\sigma C_n} - 1 \vs{0.3} \\
\dsp \dfrac{1}{a'} &=&\dsp \dfrac{4 F \sin\phi \cos\phi}{\sigma C_t}+1
\end{array}
\label{eq:aap_with_correction_prop}
\end{equation}


By combining equations \eqref{eq:aap_with_correction_prop} and  \eqref{eq:lam_prop} we get a non trivial
equation with  $\phi$ angle as variable :

\begin{empheq}[box={\mymath[colback=green!20,drop lifted shadow, sharp corners]}]{equation}
\dsp 
    F_{\rm Newton}(\phi, r) = 0 =  - 1 +    \frac{1}{\lambda_r(r)} \ \frac{2 F(\phi,r) \sin2\phi + \sigma(r) C_t(\phi)}{4 F(\phi,r) \sin^2\phi - \sigma(r) C_n(\phi)} 
\label{eq:phi+tiploss_prop}
\end{empheq}

\end{itemize}




%****************************************
\subsection{Resolution, full algorithm}
%****************************************


The algorithm is :

\begin{enumerate}
\item read the data : blade geometry, airfoil characteristics ($C_\ell, C_d$)
\item discretise the blade, define the elements and panel function of  $r$ 
\item for each elements :
      \begin{enumerate}
      \item Solve $\phi(r)$ with the equation \eqref{eq:Phi_prop}, or equation \eqref{eq:phi+tiploss_prop}
      if tip loss are taken into account,
      find the local angle of attack $\alpha$ from \eqref{eq:alpha_prop}
      \item Calculate $C_t$, $C_n$  from equations \eqref{eq:CnCt_prop} then $a$ and $a'$, from equations 
      \eqref{eq:a_prop} and \eqref{eq:ap_prop}  or  \eqref{eq:aap_with_correction_prop}
      \item Calculate the distribution of $dT$ and $dQ$
      \end{enumerate}
\item Calculate the total force, total torque, powers and efficiency.
\item Plots all the desired curves and save data.
\end{enumerate}

The crucial points, which quite fast is the Newton method in step 3a).


% \begin{figure}
% \begin{center}
% \input{Figures/propeller_side.pdf_tex}
% \label{fig:propeller_side}
% \caption{Scheme of a modern propeller, side view}
% \end{center}
% \end{figure}

% \begin{figure}
% \begin{center}
% \input{Figures/forces.pdf_tex}
% \label{fig:forces}
% \caption{Scheme of the forces on a element}
% \end{center}
% \end{figure}


\begin{figure}[htb]
  \begin{center}
  \input{Figures/pitch.pdf_tex}
  \caption{Propeller pitch.}
  \label{fig:pitch_propeller}
  \end{center}
\end{figure}



  \caption{Rotor disk  of a propeller of the Airbus A400 M.}
  \label{fig:propeller_alone}
  \end{center}
\end{figure}

A example of modern aircraft propeller is shown in figure \ref{fig:propeller_alone}.

For additional details read the section \ref{sec:wt_gmtry}


%******************************
\subsection{Upstream conditions}
%******************************

As mentioned before, the upstream velocity is assumed constant: $V(r) = V_0$ in any equation.


%***********************
\subsection{Velocities}
%***********************

\begin{figure}[htb]
  \begin{center}
  \input{Figures/propeller_side_notations.pdf_tex}
  \caption{Scheme of a modern propeller, side view (left) and airfoil section with velocity and force decompositions (right)}
  \label{fig:propeller_side_notations}
  \end{center}
  \end{figure}
  
  To develop the BEM theory we have to refer to the figure \ref{fig:propeller_side_notations} which provides
  velocity and force decompositions.
  
  \Important{
  The main difference between a rotor blade section and a propeller blade section is the direction of the airfoil with respect
  to the rotation plane. The orientation on a propeller is choosen to produce positive thrust.
  }


  This orientation naturally modifies slightly the velocity and force directions and amplitudes as explained below.


  Let's consider a given airfoil section located at a distance $r$ from the axis of rotation (fig. \ref{fig:propeller_alone}). 
  The BEMT consists on analyzing the decomposition of the \textbf{apparent velocity} vector refered as $\bs{W}$ 
  which is the air velocity seen by the airfoil section and
  which allow to determine an angle of attack $\alpha$.
  
  The apparent velocity is the sum of :
  \begin{itemize}
  \item the constant upstream  velocity $\bs{V}$
  
  \item the tangential velocity related to the disk rotation $\bs{V_t} = - r~\Omega~\bs{e_t}$,
  where $\bs{e_t} = \bs{e_x} \times \bs{e_r}$. 
  
  \item the induced velocity $\bs{w}$ whose direction is defined by the induced angle of attack $\alpha_i$. 
  It is due to the tip vortex of the blade. Its direction is assumed orthogonal to the airfoil section axis.   
  \end{itemize}
  
  
  In addition we can define the relative velocity $\bs{V_R} = \bs{V} + \bs{V_t}$ whose direction is the angle $\varphi$.

  It can be easily demonstrated that :
  \begin{equation}
  \tan \varphi =  \frac{V}{r~\Omega} = \dfrac{J}{\pi \tilde r}, \qquad \tilde r = \dfrac{2~r}{D} 
\label{eq:varphi_prop}  
\end{equation}
  
  And finally, the apparent velocity can be projected into the $\bs{e_x}$ and $\bs{e_t}$ directions as :
  
  \begin{equation}
  \bs{W} = U_x~\bs{e_x} + U_t~\bs{e_t}
  \label{eq:W_def2}
  \end{equation} 
  
  where :
  \begin{itemize}
    \item $U_x(r)$ is the axial velocity, air velocity crossing the propeller plane
    \item $U_t(r)$ is the tangential velocity, in the propeller plane. 
  \end{itemize}
  
  All these velocities are drawn on figure  \ref{fig:propeller_side_notations}.

  \Note{
  In addition the tangential induced velocity $u_t = V_t - U_t$ is drawn in blue. Because of the direction
  of the airfoil section this velocity is opposite comparing to the rotor case, and because of the direction
  of the induced velocity $w$ we can notice that $U_x > V$.}

%********************
\subsection{Angles}
%********************

Angles play an important role to define velocity and force directions.

Different angles are defined from the velocities and aerodynamics :
\begin{itemize}
  \item the induced angle of attack $\alpha_i$
  \item the relative velocity angle $\varphi$
  \item the local flow angle $\phi$ with 
  $$
  \phi = \alpha_i  + \varphi
  $$
  \item the pitching angle $\theta$, which comes from the geometry
  \item the real angle of attack $\alpha$
  \item the aerodynamic efficiency angle $\gamma$ with $\tan \gamma = C_d/C_\ell = 1/f$.
\end{itemize}

$f$ is here the aerodynamic efficiency. $C_d$ and $C_{\ell}$ are the usual lift and drag aerodynamic 
coefficients  used in the next section.


To determine forces exerted on the airfoil section we need to clearly defined the angle of attack  by :

\Important{
\begin{equation}
\alpha = \theta-\phi
\label{eq:alpha_prop}
\end{equation}
}

\gris{We can see that the angle of attack equation is the opposite for a wind turbine rotor}

%************************************
\subsection{Forces on a wing section}
%************************************

On figure \ref{fig:propeller_side_notations} are drawn the aerodynamics force on a wing section.

Let write the equations for a given section, and for a single blade :

\begin{itemize}
    \item the lift 
    $$
      \dsp \bs{dL} = q~dS~C_\ell~\bs{j_a}  
    $$
    \item the drag 
    $$
      \dsp \bs{dD} = q~dS~C_d ~\bs{i_a}
    $$
    \item the aerodynamic force $\bs{dF}$
    $$\dsp
    \bs{dF} = dL~\bs{j_a} + dD~\bs{i_a} = \bs{dT} + \bs{dF_t}
    $$
    \item the normal aerodynamic force :
    $$
    dT = \bs{dF}\cdot\bs{e_x} = q~dS~C_n
    $$
    \item the tangential aerodynamic force:
    $$
    dF_t = \bs{dF}\cdot\bs{e_y} = q~dS~C_t
    $$

\end{itemize}

with 
\begin{itemize}
  \item $q=\frac{1}{2} \rho W^2$, the local dynamic pressure
  % \item $S = c(r) \times b(r) $ where $b(r)$ is the span of the local element.
  \item $C_\ell, C_d$ are the aerodynamic lift and drag coefficients dependent on the airfoil.
  \item $C_n, C_t$ are the aerodynamic normal and tangential  coefficient dependent on the $C_{\ell}$ and $C_d$.
  \item $(\bs{i_a}, \bs{j_a})$ is the local aerodynamic base ($\bs{W} = W~\bs{i_a}$) (see fig.~\ref{fig:rotor_side_notations})
  \item $ dS = c(r)~ dr$ is the area of a blade element
\end{itemize}

In addition it it easy to show that :
\begin{eqnarray}
dT &=& \cos\phi ~ dL - \sin\phi~ dD   \label{eq:dT}\\
dF_t &=& \sin\phi ~ dL + \cos\phi~ dD \label{eq:dFt} 
\end{eqnarray}

From equations \ref{eq:dT} and \ref{eq:dFt}  it is found :
\begin{equation}
\begin{array}{lcl}
C_n &=& C_\ell \cos\phi - C_d \sin \phi \\
C_t &=& C_\ell \sin\phi + C_d \cos \phi
\end{array}
\label{eq:CnCt}
\end{equation}


%********************************************
\subsection{BEM theory algorithm}
%********************************************

To develop the Blade Element Momentum Theory, several steps are necessary and some reduced quantities are 
introduced (ever discussed in the disk rotor theory).

%******************************************************************
\subsubsection{\sc Step 1 : induced factor and apparent velocity}
%*****************************************************************
The BEMT use different reduced quantities and the most important are the induced flow factor $a$ and $a'$ defined as :

\begin{itemize}
\item the induced axial velocity factor $a$ is given by :
\Def{induced axial velocity factor $a$}{
  \begin{equation}
  U_x~/~V = 1 + a
  \label{eq:a_def_prop}
  \end{equation}
}

From the figure \ref{fig:propeller_side_notations} it can be seen that $U_x > V$, that means that $0 < a $ and $a$ should be small.


\item the induced radial velocity factor $a'$ is given by 
\Def{induced radial velocity factor $a'$}{  
\begin{equation}
    U_t  = (1-a')~V_t = (1-a') ~r~\Omega
  \label{eq:aprime_def_prop}
  \end{equation}
}



\item We can then find :
\begin{equation} \dsp
\tan \phi = \frac{U_x}{U_t} = \frac{1 + a}{1-a'} \times \frac{V}{r~\Omega}, \qquad U_x = W \sin\phi, \qquad U_t = W \cos\phi
\label{eq:tanphi_prop}
\end{equation}
\end{itemize}

It is now possible to express the modulus of the apparent velocity as a function of the induced factors :
\begin{equation}
\dsp W = \sqrt{V^2(1+a)^2 + r^2 \Omega^2 (1-a')^2} = V  \sqrt{(1+a)^2 + \left(\frac{r \Omega}{V}\right)^2 (1-a')^2}
\label{eq:W_prop}
\end{equation}


From the previous equation, we can defined a reduced velocity parameter  refers as \textbf{local tip speed ratio}
and ever seen in equation \eqref{eq:varphi_prop}:
  
\begin{equation} \dsp
  \lambda_r =  \frac{r \Omega}{V} = \frac{1}{\tan \varphi}  = \frac{\pi~\tilde r}{J}
  \label{eq:lambda_r_prop}
\end{equation}


Finally, by using the previous equations, we can show the result:

\begin{equation} \dsp
\dfrac{a'}{a} = \frac{C_t}{\lambda_r C_n}
  \label{eq:aprime_over_a}
\end{equation}



%******************************************************
\subsubsection{\sc Step 2 : equations for induction factors}
%******************************************************

The next step is to make the relationship between velocity, induced factors 
and aerodynamic coefficients.

\begin{itemize}

\item  \textbf{axial induction factor} equation.

First we can write the elementary axial force   for $N_b$ blades as :
$$ \dsp
dT = N_b~q~dS~C_n = \frac{1}{2} \rho W^2 C_n  N_b~c(r)~dr, \quad {\rm with}\qquad U_x = W \sin\phi 
$$

The  $dT$ is also given from the rotor disk theory by the
equation \eqref{eq:Tf-r2}  

Equating the both  equations leads to :
\begin{equation}
\frac{C_n~N_b~c}{8 \pi r ~\sin^2\phi} =  \frac{a}{1+a}  
\label{eq:tmp1_prop}
\end{equation}


Here can be introduced the
\textbf{solidity function}  $\sigma$ (sometimes refered as $s$): 

\Def{Solidity $\sigma$}{
\begin{equation}
  \sigma = \dfrac{N_b~c(r)}{2\pi~r} = \dfrac{N_b~c(r)}{\pi~D} ~\frac{1}{\tilde r}
\end{equation}
}
where $c(r)$ is the chord law and $N_b$ the number of blades.


For propellers, as the one shown on figure   \ref{fig:rotor_and_propeller} the solidity is much larger, but
decreases with increasing $r$.



Finally the equation \eqref{eq:tmp1_prop} can be simplified 
by introducing the solidity parameter as  

\begin{equation} \dsp \dfrac{a}{1+a} = \dfrac{\sigma C_n}{4 \sin^2\phi}, \qquad
  \dfrac{1}{a} =  \dfrac{4 \sin^2\phi}{\sigma C_n}  - 1
  \label{eq:a_prop}\end{equation} 

\item  \textbf{tangential induction factor} equation :


The elementary torque on $N_b$ blades,
 from figure \ref{fig:rotor_side_notations} is given by 
$$
dQ = N_b~r~dF_t = N_b~r~q~c(r)~dr~C_t = \frac{\rho}{2} W^2~N_b~r~c(r)~dr~C_t
$$

In this equation we rewrite 
$\dsp W^2 = W\times W  =  \frac{U_x~U_t}{\cos\phi~\sin\phi}$ with
$U_t = r~\Omega~(1-a')$

It is equal to the elementary momentum given by the rotor disk theory equation
\eqref{eq:def_dM2} written as :
$$
dQ =  4 \  \pi \ \rho \ V \ \Omega \  a' \ (1+a)  \ r^3~dr= 4~\pi~\rho~r^3~\Omega~a'~U_x~dr
$$


finally we find :


\begin{equation} \dsp 
  \dfrac{a'}{1-a'}  =   \dfrac{\sigma C_t}{4 \sin\phi \cos\phi} \qquad
  \Longrightarrow \qquad \dfrac{1}{a'} = \dfrac{4 \sin\phi \cos\phi}{\sigma C_t}+1 
\label{eq:ap_prop}\end{equation}  

\item Local flow angle $\phi$, from the equation \eqref{eq:tanphi} we get
$$\sin\phi = \dfrac{V}{W} (1+a), \qquad  \cos\phi = \dfrac{r \Omega}{W} (1-a')
$$
\begin{equation}\lambda_r \tan \phi = \dfrac{1+a}{1-a'} \label{eq:lam_prop}\end{equation}  

\gris{This equation has ever been introduced in equation \eqref{eq:tanphi_prop}.}
\end{itemize}

A single equation without the induction factors $a$ and $a'$ can be obtained from 
equations \eqref{eq:a_prop}, \eqref{eq:ap_prop} and  \eqref{eq:lam_prop}

\Important{
  \begin{equation} \dsp 
      F_{\rm Newton}(\phi, r) = 0 =  - 1 +   \frac{1}{\lambda_r(r)} \ \frac{2 \sin2\phi +\sigma(r) C_t(\phi)}{4 \sin^2\phi - \sigma(r) C_n(\phi)} 
  \label{eq:Phi_prop}
  \end{equation} 
}

%******************************************************
\subsubsection{\sc Step 3 : thrust and torque}
%******************************************************

\begin{itemize}
  \item the elementary thust is given from previous section as :
  $$ \dsp
   \der{T}{r} =  \dfrac{1}{2}~\rho~N_b~V^2~c(r)~\dfrac{(1+a)^2}{\sin^2\phi}~C_n~=~
   4 \ \pi \ \rho \ V^2 \ a \ (1+a) \ r
  $$
  
  and the total  thrust is :
  $$ \dsp
  T  = \int_{r_0}^{D/2} dT
  $$
  
  \item the elementary torque has ever been calculated  as well:
  $$\dsp
  \der{Q}{r} =  \dfrac{1}{2}~\rho~N_b~(r~\Omega)^2~r~c(r)~\dfrac{(1-a')^2}{\cos^2\phi}~C_t~=~
   4~\pi~\rho~r^3~V~\Omega~a'(1+a)
  $$
  and the total torque is 
  
  $$ \dsp
  Q  = \int_{r_0}^{D/2} dQ
  $$
  \end{itemize}

%******************************************************
\subsubsection{\sc Step 4 : power and efficiency}
%******************************************************

\begin{itemize}
  \item the input propeller power is :
  $$ \dsp
  P_{in}  = \int_{r_0}^{D/2} \Omega\times dQ = \Omega \times Q 
  $$
  
  
  \item the output thrust power is 
    
  $$ \dsp
  P_{in}   = \int_{r_0}^{D/2} V \times dT = V \times T
  $$
  
  \item the efficiency is therefore:
  $$\dsp
  \eta = \frac{P_{out}}{P_{in}} = \frac{V}{\Omega} \times \frac{T}{Q} = \frac{V^2}{\Omega^2} \times
  \frac{\dsp \int_{r_0}^{D/2}  a (1+a) r~dr }{\dsp \int_{r_0}^{D/2} a'(1+a)r^3~dr} =
  \frac{J^2}{\pi^2} \times
  \frac{\dsp \int_{\tilde r_0}^{1}  a (1+a) \tilde r~d\tilde r }{\dsp \int_{\tilde r_0}^{1} a'(1+a)\tilde r^3~d\tilde r}
  $$
  \end{itemize}
  


%*******************************************************
\subsection{Tip loss}
%*******************************************************

In the previous theory the tip loss existing because the blade is a wing of finite span has not been taken
into account. Prandtl has proposed a correction  to the BEMT theory by introducing a function $F(r)$ as follows :

\begin{itemize}
\item for elementary torque and thrust :
$$ \dsp
\der{Q}{r} = 4 \  \pi \ \rho \ V \ \Omega \  a' \ (1+a)  \ r^3~F(\phi,r)
$$

$$ \dsp
\der{T}{r} =    4 \ \pi \ \rho \ V^2 \ a \ (1+a) \ r~F(\phi,r)
$$

with
$$ \dsp
F(\phi,r) = \dfrac{2}{\pi} \arccos {\rm e}^{-f} , \qquad f =  \dfrac{N_b}{2} \ \dfrac{R-r}{r \sin\phi})
$$

\item the induction factor equations are modified  as :

\begin{equation}
\begin{array}{lcl}
\dsp \dfrac{1}{a} &=& \dsp \dfrac{4 F \sin^2\phi}{\sigma C_n} - 1 \vs{0.3} \\
\dsp \dfrac{1}{a'} &=&\dsp \dfrac{4 F \sin\phi \cos\phi}{\sigma C_t}+1
\end{array}
\label{eq:aap_with_correction_prop}
\end{equation}


By combining equations \eqref{eq:aap_with_correction_prop} and  \eqref{eq:lam_prop} we get a non trivial
equation with  $\phi$ angle as variable :

\begin{empheq}[box={\mymath[colback=green!20,drop lifted shadow, sharp corners]}]{equation}
\dsp 
    F_{\rm Newton}(\phi, r) = 0 =  - 1 +    \frac{1}{\lambda_r(r)} \ \frac{2 F(\phi,r) \sin2\phi + \sigma(r) C_t(\phi)}{4 F(\phi,r) \sin^2\phi - \sigma(r) C_n(\phi)} 
\label{eq:phi+tiploss_prop}
\end{empheq}

\end{itemize}




%****************************************
\subsection{Resolution, full algorithm}
%****************************************


The algorithm is :

\begin{enumerate}
\item read the data : blade geometry, airfoil characteristics ($C_\ell, C_d$)
\item discretise the blade, define the elements and panel function of  $r$ 
\item for each elements :
      \begin{enumerate}
      \item Solve $\phi(r)$ with the equation \eqref{eq:Phi_prop}, or equation \eqref{eq:phi+tiploss_prop}
      if tip loss are taken into account,
      find the local angle of attack $\alpha$ from \eqref{eq:alpha_prop}
      \item Calculate $C_t$, $C_n$  from equations \eqref{eq:CnCt_prop} then $a$ and $a'$, from equations 
      \eqref{eq:a_prop} and \eqref{eq:ap_prop}  or  \eqref{eq:aap_with_correction_prop}
      \item Calculate the distribution of $dT$ and $dQ$
      \end{enumerate}
\item Calculate the total force, total torque, powers and efficiency.
\item Plots all the desired curves and save data.
\end{enumerate}

The crucial points, which quite fast is the Newton method in step 3a).


% \begin{figure}
% \begin{center}
% \input{Figures/propeller_side.pdf_tex}
% \label{fig:propeller_side}
% \caption{Scheme of a modern propeller, side view}
% \end{center}
% \end{figure}

% \begin{figure}
% \begin{center}
% \input{Figures/forces.pdf_tex}
% \label{fig:forces}
% \caption{Scheme of the forces on a element}
% \end{center}
% \end{figure}


\begin{figure}[htb]
  \begin{center}
  \input{Figures/pitch.pdf_tex}
  \caption{Propeller pitch.}
  \label{fig:pitch_propeller}
  \end{center}
\end{figure}



  \caption{Rotor disk  of a propeller of the Airbus A400 M.}
  \label{fig:propeller_alone}
  \end{center}
\end{figure}

A example of modern aircraft propeller is shown in figure \ref{fig:propeller_alone}.

For additional details read the section \ref{sec:wt_gmtry}


%******************************
\subsection{Upstream conditions}
%******************************

As mentioned before, the upstream velocity is assumed constant: $V(r) = V_0$ in any equation.


%***********************
\subsection{Velocities}
%***********************

\begin{figure}[htb]
  \begin{center}
  \input{Figures/propeller_side_notations.pdf_tex}
  \caption{Scheme of a modern propeller, side view (left) and airfoil section with velocity and force decompositions (right)}
  \label{fig:propeller_side_notations}
  \end{center}
  \end{figure}
  
  To develop the BEM theory we have to refer to the figure \ref{fig:propeller_side_notations} which provides
  velocity and force decompositions.
  
  \Important{
  The main difference between a rotor blade section and a propeller blade section is the direction of the airfoil with respect
  to the rotation plane. The orientation on a propeller is choosen to produce positive thrust.
  }


  This orientation naturally modifies slightly the velocity and force directions and amplitudes as explained below.


  Let's consider a given airfoil section located at a distance $r$ from the axis of rotation (fig. \ref{fig:propeller_alone}). 
  The BEMT consists on analyzing the decomposition of the \textbf{apparent velocity} vector refered as $\bs{W}$ 
  which is the air velocity seen by the airfoil section and
  which allow to determine an angle of attack $\alpha$.
  
  The apparent velocity is the sum of :
  \begin{itemize}
  \item the constant upstream  velocity $\bs{V}$
  
  \item the tangential velocity related to the disk rotation $\bs{V_t} = - r~\Omega~\bs{e_t}$,
  where $\bs{e_t} = \bs{e_x} \times \bs{e_r}$. 
  
  \item the induced velocity $\bs{w}$ whose direction is defined by the induced angle of attack $\alpha_i$. 
  It is due to the tip vortex of the blade. Its direction is assumed orthogonal to the airfoil section axis.   
  \end{itemize}
  
  
  In addition we can define the relative velocity $\bs{V_R} = \bs{V} + \bs{V_t}$ whose direction is the angle $\varphi$.

  It can be easily demonstrated that :
  \begin{equation}
  \tan \varphi =  \frac{V}{r~\Omega} = \dfrac{J}{\pi \tilde r}, \qquad \tilde r = \dfrac{2~r}{D} 
\label{eq:varphi_prop}  
\end{equation}
  
  And finally, the apparent velocity can be projected into the $\bs{e_x}$ and $\bs{e_t}$ directions as :
  
  \begin{equation}
  \bs{W} = U_x~\bs{e_x} + U_t~\bs{e_t}
  \label{eq:W_def2}
  \end{equation} 
  
  where :
  \begin{itemize}
    \item $U_x(r)$ is the axial velocity, air velocity crossing the propeller plane
    \item $U_t(r)$ is the tangential velocity, in the propeller plane. 
  \end{itemize}
  
  All these velocities are drawn on figure  \ref{fig:propeller_side_notations}.

  \Note{
  In addition the tangential induced velocity $u_t = V_t - U_t$ is drawn in blue. Because of the direction
  of the airfoil section this velocity is opposite comparing to the rotor case, and because of the direction
  of the induced velocity $w$ we can notice that $U_x > V$.}

%********************
\subsection{Angles}
%********************

Angles play an important role to define velocity and force directions.

Different angles are defined from the velocities and aerodynamics :
\begin{itemize}
  \item the induced angle of attack $\alpha_i$
  \item the relative velocity angle $\varphi$
  \item the local flow angle $\phi$ with 
  $$
  \phi = \alpha_i  + \varphi
  $$
  \item the pitching angle $\theta$, which comes from the geometry
  \item the real angle of attack $\alpha$
  \item the aerodynamic efficiency angle $\gamma$ with $\tan \gamma = C_d/C_\ell = 1/f$.
\end{itemize}

$f$ is here the aerodynamic efficiency. $C_d$ and $C_{\ell}$ are the usual lift and drag aerodynamic 
coefficients  used in the next section.


To determine forces exerted on the airfoil section we need to clearly defined the angle of attack  by :

\Important{
\begin{equation}
\alpha = \theta-\phi
\label{eq:alpha_prop}
\end{equation}
}

\gris{We can see that the angle of attack equation is the opposite for a wind turbine rotor}

%************************************
\subsection{Forces on a wing section}
%************************************

On figure \ref{fig:propeller_side_notations} are drawn the aerodynamics force on a wing section.

Let write the equations for a given section, and for a single blade :

\begin{itemize}
    \item the lift 
    $$
      \dsp \bs{dL} = q~dS~C_\ell~\bs{j_a}  
    $$
    \item the drag 
    $$
      \dsp \bs{dD} = q~dS~C_d ~\bs{i_a}
    $$
    \item the aerodynamic force $\bs{dF}$
    $$\dsp
    \bs{dF} = dL~\bs{j_a} + dD~\bs{i_a} = \bs{dT} + \bs{dF_t}
    $$
    \item the normal aerodynamic force :
    $$
    dT = \bs{dF}\cdot\bs{e_x} = q~dS~C_n
    $$
    \item the tangential aerodynamic force:
    $$
    dF_t = \bs{dF}\cdot\bs{e_y} = q~dS~C_t
    $$

\end{itemize}

with 
\begin{itemize}
  \item $q=\frac{1}{2} \rho W^2$, the local dynamic pressure
  % \item $S = c(r) \times b(r) $ where $b(r)$ is the span of the local element.
  \item $C_\ell, C_d$ are the aerodynamic lift and drag coefficients dependent on the airfoil.
  \item $C_n, C_t$ are the aerodynamic normal and tangential  coefficient dependent on the $C_{\ell}$ and $C_d$.
  \item $(\bs{i_a}, \bs{j_a})$ is the local aerodynamic base ($\bs{W} = W~\bs{i_a}$) (see fig.~\ref{fig:rotor_side_notations})
  \item $ dS = c(r)~ dr$ is the area of a blade element
\end{itemize}

In addition it it easy to show that :
\begin{eqnarray}
dT &=& \cos\phi ~ dL - \sin\phi~ dD   \label{eq:dT}\\
dF_t &=& \sin\phi ~ dL + \cos\phi~ dD \label{eq:dFt} 
\end{eqnarray}

From equations \ref{eq:dT} and \ref{eq:dFt}  it is found :
\begin{equation}
\begin{array}{lcl}
C_n &=& C_\ell \cos\phi - C_d \sin \phi \\
C_t &=& C_\ell \sin\phi + C_d \cos \phi
\end{array}
\label{eq:CnCt}
\end{equation}


%********************************************
\subsection{BEM theory algorithm}
%********************************************

To develop the Blade Element Momentum Theory, several steps are necessary and some reduced quantities are 
introduced (ever discussed in the disk rotor theory).

%******************************************************************
\subsubsection{\sc Step 1 : induced factor and apparent velocity}
%*****************************************************************
The BEMT use different reduced quantities and the most important are the induced flow factor $a$ and $a'$ defined as :

\begin{itemize}
\item the induced axial velocity factor $a$ is given by :
\Def{induced axial velocity factor $a$}{
  \begin{equation}
  U_x~/~V = 1 + a
  \label{eq:a_def_prop}
  \end{equation}
}

From the figure \ref{fig:propeller_side_notations} it can be seen that $U_x > V$, that means that $0 < a $ and $a$ should be small.


\item the induced radial velocity factor $a'$ is given by 
\Def{induced radial velocity factor $a'$}{  
\begin{equation}
    U_t  = (1-a')~V_t = (1-a') ~r~\Omega
  \label{eq:aprime_def_prop}
  \end{equation}
}



\item We can then find :
\begin{equation} \dsp
\tan \phi = \frac{U_x}{U_t} = \frac{1 + a}{1-a'} \times \frac{V}{r~\Omega}, \qquad U_x = W \sin\phi, \qquad U_t = W \cos\phi
\label{eq:tanphi_prop}
\end{equation}
\end{itemize}

It is now possible to express the modulus of the apparent velocity as a function of the induced factors :
\begin{equation}
\dsp W = \sqrt{V^2(1+a)^2 + r^2 \Omega^2 (1-a')^2} = V  \sqrt{(1+a)^2 + \left(\frac{r \Omega}{V}\right)^2 (1-a')^2}
\label{eq:W_prop}
\end{equation}


From the previous equation, we can defined a reduced velocity parameter  refers as \textbf{local tip speed ratio}
and ever seen in equation \eqref{eq:varphi_prop}:
  
\begin{equation} \dsp
  \lambda_r =  \frac{r \Omega}{V} = \frac{1}{\tan \varphi}  = \frac{\pi~\tilde r}{J}
  \label{eq:lambda_r_prop}
\end{equation}


Finally, by using the previous equations, we can show the result:

\begin{equation} \dsp
\dfrac{a'}{a} = \frac{C_t}{\lambda_r C_n}
  \label{eq:aprime_over_a}
\end{equation}



%******************************************************
\subsubsection{\sc Step 2 : equations for induction factors}
%******************************************************

The next step is to make the relationship between velocity, induced factors 
and aerodynamic coefficients.

\begin{itemize}

\item  \textbf{axial induction factor} equation.

First we can write the elementary axial force   for $N_b$ blades as :
$$ \dsp
dT = N_b~q~dS~C_n = \frac{1}{2} \rho W^2 C_n  N_b~c(r)~dr, \quad {\rm with}\qquad U_x = W \sin\phi 
$$

The  $dT$ is also given from the rotor disk theory by the
equation \eqref{eq:Tf-r2}  

Equating the both  equations leads to :
\begin{equation}
\frac{C_n~N_b~c}{8 \pi r ~\sin^2\phi} =  \frac{a}{1+a}  
\label{eq:tmp1_prop}
\end{equation}


Here can be introduced the
\textbf{solidity function}  $\sigma$ (sometimes refered as $s$): 

\Def{Solidity $\sigma$}{
\begin{equation}
  \sigma = \dfrac{N_b~c(r)}{2\pi~r} = \dfrac{N_b~c(r)}{\pi~D} ~\frac{1}{\tilde r}
\end{equation}
}
where $c(r)$ is the chord law and $N_b$ the number of blades.


For propellers, as the one shown on figure   \ref{fig:rotor_and_propeller} the solidity is much larger, but
decreases with increasing $r$.



Finally the equation \eqref{eq:tmp1_prop} can be simplified 
by introducing the solidity parameter as  

\begin{equation} \dsp \dfrac{a}{1+a} = \dfrac{\sigma C_n}{4 \sin^2\phi}, \qquad
  \dfrac{1}{a} =  \dfrac{4 \sin^2\phi}{\sigma C_n}  - 1
  \label{eq:a_prop}\end{equation} 

\item  \textbf{tangential induction factor} equation :


The elementary torque on $N_b$ blades,
 from figure \ref{fig:rotor_side_notations} is given by 
$$
dQ = N_b~r~dF_t = N_b~r~q~c(r)~dr~C_t = \frac{\rho}{2} W^2~N_b~r~c(r)~dr~C_t
$$

In this equation we rewrite 
$\dsp W^2 = W\times W  =  \frac{U_x~U_t}{\cos\phi~\sin\phi}$ with
$U_t = r~\Omega~(1-a')$

It is equal to the elementary momentum given by the rotor disk theory equation
\eqref{eq:def_dM2} written as :
$$
dQ =  4 \  \pi \ \rho \ V \ \Omega \  a' \ (1+a)  \ r^3~dr= 4~\pi~\rho~r^3~\Omega~a'~U_x~dr
$$


finally we find :


\begin{equation} \dsp 
  \dfrac{a'}{1-a'}  =   \dfrac{\sigma C_t}{4 \sin\phi \cos\phi} \qquad
  \Longrightarrow \qquad \dfrac{1}{a'} = \dfrac{4 \sin\phi \cos\phi}{\sigma C_t}+1 
\label{eq:ap_prop}\end{equation}  

\item Local flow angle $\phi$, from the equation \eqref{eq:tanphi} we get
$$\sin\phi = \dfrac{V}{W} (1+a), \qquad  \cos\phi = \dfrac{r \Omega}{W} (1-a')
$$
\begin{equation}\lambda_r \tan \phi = \dfrac{1+a}{1-a'} \label{eq:lam_prop}\end{equation}  

\gris{This equation has ever been introduced in equation \eqref{eq:tanphi_prop}.}
\end{itemize}

A single equation without the induction factors $a$ and $a'$ can be obtained from 
equations \eqref{eq:a_prop}, \eqref{eq:ap_prop} and  \eqref{eq:lam_prop}

\Important{
  \begin{equation} \dsp 
      F_{\rm Newton}(\phi, r) = 0 =  - 1 +   \frac{1}{\lambda_r(r)} \ \frac{2 \sin2\phi +\sigma(r) C_t(\phi)}{4 \sin^2\phi - \sigma(r) C_n(\phi)} 
  \label{eq:Phi_prop}
  \end{equation} 
}

%******************************************************
\subsubsection{\sc Step 3 : thrust and torque}
%******************************************************

\begin{itemize}
  \item the elementary thust is given from previous section as :
  $$ \dsp
   \der{T}{r} =  \dfrac{1}{2}~\rho~N_b~V^2~c(r)~\dfrac{(1+a)^2}{\sin^2\phi}~C_n~=~
   4 \ \pi \ \rho \ V^2 \ a \ (1+a) \ r
  $$
  
  and the total  thrust is :
  $$ \dsp
  T  = \int_{r_0}^{D/2} dT
  $$
  
  \item the elementary torque has ever been calculated  as well:
  $$\dsp
  \der{Q}{r} =  \dfrac{1}{2}~\rho~N_b~(r~\Omega)^2~r~c(r)~\dfrac{(1-a')^2}{\cos^2\phi}~C_t~=~
   4~\pi~\rho~r^3~V~\Omega~a'(1+a)
  $$
  and the total torque is 
  
  $$ \dsp
  Q  = \int_{r_0}^{D/2} dQ
  $$
  \end{itemize}

%******************************************************
\subsubsection{\sc Step 4 : power and efficiency}
%******************************************************

\begin{itemize}
  \item the input propeller power is :
  $$ \dsp
  P_{in}  = \int_{r_0}^{D/2} \Omega\times dQ = \Omega \times Q 
  $$
  
  
  \item the output thrust power is 
    
  $$ \dsp
  P_{in}   = \int_{r_0}^{D/2} V \times dT = V \times T
  $$
  
  \item the efficiency is therefore:
  $$\dsp
  \eta = \frac{P_{out}}{P_{in}} = \frac{V}{\Omega} \times \frac{T}{Q} = \frac{V^2}{\Omega^2} \times
  \frac{\dsp \int_{r_0}^{D/2}  a (1+a) r~dr }{\dsp \int_{r_0}^{D/2} a'(1+a)r^3~dr} =
  \frac{J^2}{\pi^2} \times
  \frac{\dsp \int_{\tilde r_0}^{1}  a (1+a) \tilde r~d\tilde r }{\dsp \int_{\tilde r_0}^{1} a'(1+a)\tilde r^3~d\tilde r}
  $$
  \end{itemize}
  


%*******************************************************
\subsection{Tip loss}
%*******************************************************

In the previous theory the tip loss existing because the blade is a wing of finite span has not been taken
into account. Prandtl has proposed a correction  to the BEMT theory by introducing a function $F(r)$ as follows :

\begin{itemize}
\item for elementary torque and thrust :
$$ \dsp
\der{Q}{r} = 4 \  \pi \ \rho \ V \ \Omega \  a' \ (1+a)  \ r^3~F(\phi,r)
$$

$$ \dsp
\der{T}{r} =    4 \ \pi \ \rho \ V^2 \ a \ (1+a) \ r~F(\phi,r)
$$

with
$$ \dsp
F(\phi,r) = \dfrac{2}{\pi} \arccos {\rm e}^{-f} , \qquad f =  \dfrac{N_b}{2} \ \dfrac{R-r}{r \sin\phi})
$$

\item the induction factor equations are modified  as :

\begin{equation}
\begin{array}{lcl}
\dsp \dfrac{1}{a} &=& \dsp \dfrac{4 F \sin^2\phi}{\sigma C_n} - 1 \vs{0.3} \\
\dsp \dfrac{1}{a'} &=&\dsp \dfrac{4 F \sin\phi \cos\phi}{\sigma C_t}+1
\end{array}
\label{eq:aap_with_correction_prop}
\end{equation}


By combining equations \eqref{eq:aap_with_correction_prop} and  \eqref{eq:lam_prop} we get a non trivial
equation with  $\phi$ angle as variable :

\begin{empheq}[box={\mymath[colback=green!20,drop lifted shadow, sharp corners]}]{equation}
\dsp 
    F_{\rm Newton}(\phi, r) = 0 =  - 1 +    \frac{1}{\lambda_r(r)} \ \frac{2 F(\phi,r) \sin2\phi + \sigma(r) C_t(\phi)}{4 F(\phi,r) \sin^2\phi - \sigma(r) C_n(\phi)} 
\label{eq:phi+tiploss_prop}
\end{empheq}

\end{itemize}




%****************************************
\subsection{Resolution, full algorithm}
%****************************************


The algorithm is :

\begin{enumerate}
\item read the data : blade geometry, airfoil characteristics ($C_\ell, C_d$)
\item discretise the blade, define the elements and panel function of  $r$ 
\item for each elements :
      \begin{enumerate}
      \item Solve $\phi(r)$ with the equation \eqref{eq:Phi_prop}, or equation \eqref{eq:phi+tiploss_prop}
      if tip loss are taken into account,
      find the local angle of attack $\alpha$ from \eqref{eq:alpha_prop}
      \item Calculate $C_t$, $C_n$  from equations \eqref{eq:CnCt_prop} then $a$ and $a'$, from equations 
      \eqref{eq:a_prop} and \eqref{eq:ap_prop}  or  \eqref{eq:aap_with_correction_prop}
      \item Calculate the distribution of $dT$ and $dQ$
      \end{enumerate}
\item Calculate the total force, total torque, powers and efficiency.
\item Plots all the desired curves and save data.
\end{enumerate}

The crucial points, which quite fast is the Newton method in step 3a).


% \begin{figure}
% \begin{center}
% \input{Figures/propeller_side.pdf_tex}
% \label{fig:propeller_side}
% \caption{Scheme of a modern propeller, side view}
% \end{center}
% \end{figure}

% \begin{figure}
% \begin{center}
% \input{Figures/forces.pdf_tex}
% \label{fig:forces}
% \caption{Scheme of the forces on a element}
% \end{center}
% \end{figure}


\begin{figure}[htb]
  \begin{center}
  \input{Figures/pitch.pdf_tex}
  \caption{Propeller pitch.}
  \label{fig:pitch_propeller}
  \end{center}
\end{figure}



% \caption{Rotor disk  of a modern propeller.}
% \label{fig:propeller}
% \end{center}
% \end{figure}



% \begin{figure}
% \begin{center}
% \input{Figures/propeller_side.pdf_tex}
% \label{fig:propeller_side}
% \caption{Scheme of a modern propeller, side view}
% \end{center}
% \end{figure}

% \begin{figure}
% \begin{center}
% \input{Figures/forces.pdf_tex}
% \label{fig:forces}
% \caption{Scheme of the forces on a element}
% \end{center}
% \end{figure}

%******************************
\subsection{Upstream conditions}
%******************************

The wind turbine  is submitted to an upstream flow. The basic assumption is to consider it
uniform at velocity  $V_0$.

However, for large rotor size of wind turbine, this assumption is no longer true, because of the 
ground boundary layer. In this situation the upstream velocity is a profile $V(r)$ with some shear.
The velocity $V_0$ is then defined as a mean value given by 

$$ \dsp
V_0 = \frac{1}{S} \int_{r_0}^R V(r) dS
$$

where $S = \pi~ (R^2-r_0^2)$ is the wetted (or useful) rotor disk area.

In the following the upstream velocity is therefore refers as $V(r)$ or simply $V$.


%***********************
\subsection{Velocities}
%***********************

To develop the BEM theory we have to refer to the figure \ref{fig:rotor_side_notations} which provides
velocity and force decompositions.



Let's consider a given airfoil section located at a distance $r$ from the axis of rotation (fig. \ref{fig:blades}). 
The BEMT consists on analyzing the decomposition of the \textbf{apparent velocity} vector refered as $\bs{W}$ 
which is the air velocity seen by the airfoil section and
which allow to determine an angle of attack $\alpha$.

The apparent velocity is the sum of :
\begin{itemize}
\item the upstream  velocity $\bs{V}$

\item the tangential velocity related to the disk rotation $\bs{V_t} = - r~\Omega~\bs{e_t}$,
 where $\bs{e_t} = \bs{e_x} \times \bs{e_r}$. 

\item the induced velocity $\bs{w}$ whose direction is defined by the induced angle of attack $\alpha_i$. 
It is due to the tip vortex of the blade. Its direction is assumed orthogonal to the airfoil section axis.   
\end{itemize}


In addition we can define the relative velocity $\bs{V_R} = \bs{V} + \bs{V_t}$ whose direction is the angle $\varphi$, with 

$$ \dsp \tan\varphi = \frac{V(r)}{|V_t|} = \frac{V(r)}{r\Omega}$$

And finally, the apparent velocity can be projected into the $\bs{e_x}$ and $\bs{e_t}$ directions as :

\begin{equation}
 \bs{W} = U_x~\bs{e_x} + U_t~\bs{e_t}
\label{eq:W_def}
\end{equation} 

where :
\begin{itemize}
  \item $U_x(r)$ is the axial velocity, air velocity crossing the rotor plane
  \item $U_t(r)$ is the tangential velocity, in the rotor plane. 
\end{itemize}

All these velocities are drawn on figure  \ref{fig:rotor_side_notations}.

\begin{figure}[htb]
  \begin{center}
  \input{Figures/rotor_side_notations.pdf_tex}
  \caption{Scheme of a wind turbine rotor, side view (left) and airfoil section with velocity and force decompositions (right)}
  \label{fig:rotor_side_notations}
  \end{center}
  \end{figure}


In addition on the close-up of fig.  \ref{fig:rotor_side_notations} is shown the tangential induced velocity $u_t = U_t - r~\Omega$
discussed in section \ref{secsec:ut}.

%********************
\subsection{Angles}
%********************

Angles play an important role to define velocity and force directions.

Different angles are defined from the velocities and aerodynamics :
\begin{itemize}
  \item the induced angle of attack $\alpha_i$
  \item the relative velocity angle $\varphi$
  \item the local flow angle $\phi$ with 
  $$
    \phi = \varphi - \alpha_i
  $$

  \item the pitching angle $\theta$, which comes from the geometry
  \item the real angle of attack $\alpha$
  \item the aerodynamic efficiency angle $\gamma$ with $\tan \gamma = C_d/C_\ell = 1/f$.
\end{itemize}

$f$ is here the aerodynamic efficiency. $C_d$ and $C_{\ell}$ are the usual lift and drag aerodynamic 
coefficients  used in the next section.


To determine forces exerted on the airfoil section we need to clearly defined the angle of attack  by :

\begin{equation}
\alpha =  \phi-\theta
\label{eq:alpha}
\end{equation}

%************************************
\subsection{Forces on a wing section}
%************************************

On figure \ref{fig:rotor_side_notations} are drawn the aerodynamics force on a wing section.

Let write the equations for a given section, and for a single blade :

\begin{itemize}
    \item the lift 
    $$
      \dsp \bs{dL} = q~dS~C_\ell~\bs{j_a}  
    $$
    \item the drag 
    $$
      \dsp \bs{dD} = q~dS~C_d ~\bs{i_a}
    $$
    \item the aerodynamic force $\bs{dF}$
    $$\dsp
    \bs{dF} = dL~\bs{j_a} + dD~\bs{i_a} = \bs{dT} + \bs{dF_t}
    $$
    \item the normal aerodynamic force :
    $$
    dT = \bs{dF}\cdot\bs{e_x} = q~dS~C_n
    $$
    \item the tangential aerodynamic force:
    $$
    dF_t = \bs{dF}\cdot\bs{e_y} = q~dS~C_t
    $$

\end{itemize}

with 
\begin{itemize}
  \item $q=\frac{1}{2} \rho W^2$, the local dynamic pressure
  % \item $S = c(r) \times b(r)$ where $b(r)$ is the span of the local element.
  \item $C_\ell, C_d$ are the aerodynamic lift and drag coefficients dependent on the airfoil.
  \item $C_n, C_t$ are the aerodynamic normal and tangential  coefficient dependent on the $C_{\ell}$ and $C_d$.
  \item $(\bs{i_a}, \bs{j_a})$ is the local aerodynamic base ($\bs{W} = W~\bs{i_a}$) (see fig.~\ref{fig:rotor_side_notations})
  \item $ dS = c(r)~ dr$ is the area of a blade element
\end{itemize}

In addition it it easy to show that :
\begin{eqnarray}
dT &=& \cos\phi ~ dL + \sin\phi~ dD   \label{eq:dT_prop}\\
dF_t &=& \sin\phi ~ dL - \cos\phi~ dD \label{eq:dFt_prop} 
\end{eqnarray}

From equations \ref{eq:dT_prop} and \ref{eq:dFt_prop}  it is found :
\begin{equation}
\begin{array}{lcl}
C_n &=& C_\ell \cos\phi + C_d \sin \phi \\
C_t &=& C_\ell \sin\phi - C_d \cos \phi
\end{array}
\label{eq:CnCt_prop}
\end{equation}

\Note{The sign in front of $C_d$ is opposite between turbine and propeller:
\begin{itemize}
  \item In propeller, airfoil drag has a negative effect on the thrust, which is the system output
  \item In turbine, airfoil drag as a negative effect on $dF_t$, i.e. the momentum $dQ$, which is the system output
\end{itemize}
}

%********************************************
\subsection{BEM theory algorithm}
%********************************************


To develop the Blade Element Momentum Theory, several steps are necessary and some reduced quantities are 
introduced (ever discussed in the disk rotor theory).

%******************************************************************
\subsubsection{\sc Step 1 : induced factor and apparent velocity}
%*****************************************************************


The BEMT use different reduced quantities and the most important are the induced flow factor $a$ and $a'$ defined as :

\begin{itemize}
\item the induced axial velocity factor $a$ is given by :
  \begin{equation}
  U_x~/~V = 1 - a
  \label{eq:a_def}
  \end{equation}

From the figure \ref{fig:rotor_side_notations} it can be seen that $U_x < V$, that means that $0 < a < 1$.

\item the induced radial velocity factor $a'$ is given by 
  \begin{equation}
    U_t  = (1+a')~V_t = (1+a') ~r~\Omega
  \label{eq:aprime_def}
  \end{equation}

\item We can then find :
\begin{equation} \dsp
\tan \phi = \frac{U_x}{U_t} = \frac{1 - a}{1+a'} \times \frac{V}{r~\Omega}, \qquad U_x = W \sin\phi, \qquad U_t = W \cos\phi
\label{eq:tanphi}
\end{equation}
\end{itemize}

It is now possible to express the modulus of the apparent velocity as a function of the induced factors :
\begin{equation}
\dsp W = \sqrt{V^2(1-a)^2 + r^2 \Omega^2 (1+a')^2} = V  \sqrt{(1-a)^2 + \left(\frac{r \Omega}{V}\right)^2 (1+a')^2}
\label{eq:W}
\end{equation}


From the previous equation, we can defined a reduced velocity parameter  refers as \textbf{local tip speed ratio}:

\begin{equation}
    \lambda_r =  \frac{r \Omega}{V(r)} = \frac{1}{\tan \varphi}
  \label{eq:lambda_r}
\end{equation}

It has been also encountered in the rotor disk theory.


%******************************************************
\subsubsection{\sc Step 2 : equations for induction factors}
%******************************************************

The next step is to make the relationship between velocity, induced factors 
and aerodynamic coefficients.

\begin{itemize}

\item  \textbf{axial induction factor} equation.

First we can write the elementary axial force \underline{exerted by the fluid on the element}, for $N_b$ blades as :
$$ \dsp
dT = N_b~q~dS~C_n = \frac{1}{2} \rho W^2 C_n  N_b~c(r)~dr, \quad {\rm with}\qquad U_x = W \sin\phi 
$$

The  $dT$ is also given from the rotor disk theory by the
equation \eqref{eq:Tf-r} (or the opposite of equation \eqref{eq:3})

Equating the both  equations leads to :
\begin{equation}
\frac{C_n~N_b~c}{8 \pi r ~\sin^2\phi} = \frac{V}{U_x}-1 = \frac{1}{1-a} -1 
\label{eq:tmp1}
\end{equation}


Here can be introduced the
\textbf{solidity function}  $\sigma$ (sometimes refered as $s$): 
    
\begin{equation}
  \sigma = \dfrac{N_b~c(r)}{2\pi r}
\end{equation}

where $c(r)$ is the chord law and $N_b$ the number of blades.

For a wind turbine rotor, there are few blades (two to five) and the total area of the blade
is much lower than the disk aera, therefore the solidity is a small number.


Finally the equation \eqref{eq:tmp1} can be simplified 
by introducing the solidity parameter as  

\begin{equation} \dsp \dfrac{1}{a} = 1+ \dfrac{4 \sin^2\phi}{\sigma C_n}, \qquad   
 \dfrac{a}{1-a} = \dfrac{\sigma C_n}{4 \sin^2\phi}
  \label{eq:a}\end{equation} 

\item  \textbf{tangential induction factor} :


The elementary torque on $N_b$ blades,
 from figure \ref{fig:rotor_side_notations} is given by 
$$
dQ = N_b~r~dF_t = N_b~r~q~c(r)~dr~C_t = \frac{\rho}{2} W^2~N_b~r~c(r)~dr~C_t
$$

In this equation we rewrite 
$\dsp W^2 = W\times W  =  \frac{U_x~U_t}{\cos\phi~\sin\phi}$ with
$U_t = r~\Omega~(1+a')$

It is equal to the elementary momentum given by the rotor disk theory equation
\eqref{eq:def_dM} written as :
$$
dQ =  4 \  \pi \ \rho \ V(r) \ \Omega \  a' \ (1-a)  \ r^3~dr= 4~\pi~\rho~r^3~\Omega~a'~U_x~dr
$$


finally we find :


\begin{equation} \dsp 
  \dfrac{a'}{1+a'}  =   \dfrac{\sigma C_t}{4 \sin\phi \cos\phi} \qquad
  \Longrightarrow \qquad \dfrac{1}{a'} = \dfrac{4 \sin\phi \cos\phi}{\sigma C_t}-1 
\label{eq:ap}\end{equation}  

\item Local flow angle $\phi$, from the equation \eqref{eq:tanphi} we get
$$\sin\phi = \dfrac{V}{W} (1-a), \qquad  \cos\phi = \dfrac{r \Omega}{W} (1+a')
$$
\begin{equation}\lambda_r \tan \phi = \dfrac{1-a}{1+a'} \label{eq:lam}\end{equation}  

\end{itemize}

%******************************************************
\subsubsection{\sc Step 3 : thrust and torque}
%******************************************************

\begin{itemize}
\item the elementary thust is given from previous section as :
$$ \dsp
 \der{T}{r} =  \dfrac{1}{2}~\rho~N_b~V^2~c(r)~\dfrac{(1-a)^2}{\sin^2\phi}~C_n~=~
 4 \ \pi \ \rho \ V^2(r) \ a \ (1-a) \ r
$$

and the total  thrust is :
$$ \dsp
T  = \int_{r_0}^{R} dT
$$

\item the elementary torque has ever been calculated  as well:
$$\dsp
\der{Q}{r} =  \dfrac{1}{2}~\rho~N_b~(r~\Omega)^2~r~c(r)~\dfrac{(1+a')^2}{\cos^2\phi}~C_t~=~
 4~\pi~\rho~r^3~V(r)~\Omega~a'(1-a)
$$
and the total torque is 

$$ \dsp
Q  = \int_{r_0}^{R} dQ
$$
\end{itemize}



%******************************************************
\subsubsection{\sc Step 4 : power and efficiency}
%******************************************************

\begin{itemize}
\item the input power is :
$$ \dsp
P_{in}   = \int_{r_0}^{R} V(r) \times dT
$$

\item the output power, obtained from the torque is 
  

$$ \dsp
P_{out}  = \int_{r_0}^{R} \Omega\times dQ = \Omega \times Q 
$$
\item the efficiency is therefore:
$$\dsp
\eta = \frac{P_{out}}{P_{in}} 
$$
\end{itemize}


%*******************************************************
\subsection{Tip loss}
%*******************************************************

In the previous theory the tip loss existing because the blade is a wing of finite span has not been taken
into account. Prandtl has proposed a correction  to the BEMT theory by introducing a function $F(r)$ as follows :

\begin{itemize}
\item for elementary torque and thrust :
$$ \dsp
\der{Q}{r} = 4 \  \pi \ \rho \ V_0(r) \ \Omega \  a' \ (1-a)  \ r^3~F(r)
$$

$$ \dsp
\der{T}{r} =    4 \ \pi \ \rho \ V_0^2 \ a \ (1-a) \ r~F(r)
$$

with
$$ \dsp
F(\phi,r) = \dfrac{2}{\pi} \arccos e^{-f} , \qquad f =  \dfrac{N_b}{2} \ \dfrac{R-r}{r \sin\phi}
$$

\item the induction factor equations are modified  as :

\begin{equation}
\begin{array}{lcl}
\dsp \dfrac{1}{a} &=& \dsp 1+ \dfrac{4 F \sin^2\phi}{\sigma C_n} \vs{0.3} \\
\dsp \dfrac{1}{a'} &=&\dsp \dfrac{4 F \sin\phi \cos\phi}{\sigma C_t}-1
\end{array}
\label{eq:aap_with_correction}
\end{equation}


By combining equations \eqref{eq:aap_with_correction} and  \eqref{eq:lam} we get a non trivial
equation with  $\phi$ angle as variable :

\begin{empheq}[box={\mymath[colback=green!20,drop lifted shadow, sharp corners]}]{equation}
\dsp 
    f_{\rm Newton} (\phi, r)= 0 = -1 +  \frac{1}{\lambda_r(r)} \ \frac{2 F(\phi,r) \sin2\phi - \sigma(r) C_t(\phi)}{4 F(\phi,r) \sin^2\phi + \sigma(r) C_n(\phi)} 
\label{eq:phi+tiploss}
\end{empheq}

\end{itemize}



%****************************************
\subsection{Resolution, full algorithm}
%****************************************

To solve the problem an iteration procedure is required to first determined $a$ and $a'$ for every element, and 
a basic procedure, always retains by people have been proposed by Glauert.

Here we propose something new:

By combining equations \eqref{eq:a}, \eqref{eq:ap} and  \eqref{eq:lam} we get a non trivial
equation with  $\phi$ angle as variable :

\begin{empheq}[box={\mymath[colback=green!20,drop lifted shadow, sharp corners]}]{equation}
\dsp 
f_{\rm Newton} (\phi, r)= 0 = -1   \frac{1}{\lambda_r(r)} \ \frac{2 \sin2\phi - \sigma(r) C_t(\phi)}{4 \sin^2\phi + \sigma(r) C_n(\phi)} 
\label{eq:Phi}
\end{empheq}

where $C_n$ and $C_t$ depend on  $\alpha$,  $\phi$ and of the pitching law $\theta(r)$.

This procedure only requires a Newton procedure to solve \eqref{eq:Phi}, and then all the unknowns are calculated 
from $\phi$, $a$ and $a'$.

The algorithm is now :

\begin{enumerate}
\item read the data : blade geometry, airfoil characteristics ($C_\ell, C_d$)
\item discretise the blade, define the elements and panel function of  $r$ 
\item for each elements :
      \begin{enumerate}
      \item Solve $\phi(r)$ with the equation \eqref{eq:Phi}, or equation \eqref{eq:phi+tiploss} if tip loss are taken into account,
      find the local angle of attack $\alpha$ from \eqref{eq:alpha}
      \item Calculate $C_t$, $C_n$  from equations \eqref{eq:CnCt} then $a$ and $a'$, from equations  \eqref{eq:a} and \eqref{eq:ap}  or  \eqref{eq:aap_with_correction}
      \item Calculate the distribution of $dT$ and $dQ$
      \end{enumerate}
\item Calculate the total force, total torque, powers and efficiency.
\item Plots all the desired curves and save data.
\end{enumerate}

The crucial points, which quite fast is the Newton method in step 3a).
